\chapter{Categories}
\label{cap:categories}

\begin{objectius}
\apartat{Prerequisits} Llenguatge elemental de conjunts i aplicacions: domini, codomini i composició; producte cartesià, subconjunts i relacions. Lectura de quantificadors elementals i distinció entre existència i unicitat. Nocions bàsiques d'implicació, equivalència i contraexemple. Els exemples amb grups, espais vectorials o espais topològics enriqueixen la lectura, però no són imprescindibles. No es pressuposa cap coneixement previ de teoria de categories.
\apartat{Objectius} En acabar el capítol, el lector hauria de poder:
\begin{itemize}
  \item enunciar la definició de categoria i verificar-ne els axiomes en exemples concrets;
  \item llegir i escriure diagrames commutatius i traduir-los a igualtats de composicions;
  \item reconèixer la mateixa estructura categòrica en conjunts, preordres, monoides, grafs, relacions i sistemes deductius;
  \item definir i utilitzar isomorfismes, monomorfismes, epimorfismes, seccions i retraccions, distingint-los de les nocions conjuntistes corresponents;
  \item construir i interpretar la categoria oposada, categories producte, subcategories i categories sobre i sota un objecte en casos senzills;
  \item aplicar el principi de dualitat a definicions i resultats elementals;
  \item descriure la categoria lliure generada per un graf finit senzill;
  \item distingir categories petites i localment petites.
\end{itemize}
\end{objectius}

\section{Cap al concepte de categoria}

Quan comencem a estudiar matemàtiques, aprenem a treballar amb molts tipus d'\emph{estructures}: ordres, grups, espais vectorials, espais topològics, etc. Aviat notem que aquestes estructures rarament apareixen soles. Al costat de cadascuna hi ha sempre una classe de transformacions que les relacionen respectant la seva manera de comportar-se: entre ordres, les aplicacions monòtones; entre grups, els homomorfismes; entre espais vectorials, les aplicacions lineals; entre espais topològics, les aplicacions contínues. Aquestes transformacions admissibles les anomenarem, en general, \emph{morfismes} o \emph{fletxes}.

La noció de categoria no decideix per si sola quins morfismes són admissibles: ho fixem nosaltres en cada context, d'acord amb l'estructura que volem respectar. El que fa la teoria de categories és abstreure el patró comú a tots aquests contextos: treballar amb objectes, amb transformacions admissibles entre ells i amb la manera com aquestes transformacions es componen. Cada categoria concreta enregistra aquesta tria de morfismes i, a partir d'aquí, només en reté la composició i les identitats.

Els exemples concrets que acabem d'esmentar tenen tots un conjunt subjacent. En aquests casos, preservar l'estructura significa començar per una aplicació entre els conjunts subjacents i imposar-hi les condicions adequades. No és, però, un requisit de la noció de categoria: més endavant trobarem categories els objectes de les quals no són conjunts amb estructura. Com hem dit a la secció~\ref{sec:mirar-fora} de la introducció, la teoria de categories desplaça la visió clàssica, puntual i conjuntista, de les aplicacions cap a una visió abstracta com a fletxes.

Recordem breument la notació de fletxa per a les aplicacions entre conjunts. Una aplicació $f$ amb domini $X$ i codomini $Y$ es denota $f\colon X\to Y$, o bé com a fletxa $X\xrightarrow{f}Y$. L'operació fonamental sobre les aplicacions és la \emph{composició}: si $f\colon X\to Y$ i $g\colon Y\to Z$, es defineix $g\comp f\colon X\to Z$ per
\[
(g\comp f)(x)\coloneqq g(f(x)).
\]
Perquè la composició estigui definida, el codomini de $f$ ha de coincidir amb el domini de $g$; amb fletxes, $X\xrightarrow{f}Y\xrightarrow{g}Z$. A més, per a cada conjunt $X$ hi ha l'aplicació \emph{identitat} sobre $X$,
\[
\id_X\colon X\to X,\qquad \id_X(x)\coloneqq x.
\]
Aquestes operacions es regeixen per la llei d'associativitat i les lleis de la unitat: per a $f\colon X\to Y$, $g\colon Y\to Z$ i $h\colon Z\to W$,
\[
(h\comp g)\comp f=h\comp(g\comp f),\qquad f\comp\id_X=f=\id_Y\comp f.
\]
Observem que aquestes equacions es formulen purament en termes de les operacions sobre aplicacions, sense cap referència als elements dels conjunts $X$, $Y$, $Z$, $W$.

Això no és cap ingenuïtat, sinó una evidència sobre les estructures que hem mencionat: els homomorfismes entre grups, les aplicacions monòtones entre preordres, les aplicacions lineals entre espais vectorials i les aplicacions contínues entre espais topològics satisfan les condicions que acabem d'escriure.

Aquesta és la idea que orientarà tot el que segueix: estudiar un objecte no pel que és per dins, sinó pel paper que fa dins del sistema de fletxes que el connecten amb els altres. És el que ens permetrà de parlar-ne amb un mateix llenguatge i, fins i tot, estendre'l a altres tipus d'objectes no necessàriament basats en la noció d'aplicació, com ara grafs, sistemes deductius o relacions entre conjunts.

\section{Definició de categoria}

\begin{definicio}
\label{def:categoria}
\index{categoria}Una \textbf{categoria} $\cat{C}$ consta de les dades següents:
\begin{enumerate}
\item una col·lecció d'\textbf{objectes}\index{objecte} $A,B,C,\dots$, que denotem $\Ob(\cat{C})$;
\item una col·lecció de \textbf{morfismes}\index{morfisme} $f,g,h,\dots$, que denotem $\Mor(\cat{C})$; cada morfisme $f$ té associats un únic \textbf{domini}\index{domini} $A$ i un únic \textbf{codomini}\index{codomini} $B$, i ho escrivim com $A=\dom(f)$ i $B=\cod(f)$; en aquest cas també escrivim $f\colon A\to B$ o com a fletxa $A\xrightarrow{f}B$;
\item per a cada objecte $A$, un morfisme \textbf{identitat}\index{identitat} $\id_A\colon A\to A$;
\item per a cada parella de morfismes componibles $f\colon A\to B$ i $g\colon B\to C$, la \textbf{composició}\index{composicio@composició} $g\comp f\colon A\to C$.
\end{enumerate}
Aquestes dades han de satisfer els dos axiomes següents, sempre que les composicions indicades tinguin sentit:
\begin{enumerate}
\item \textbf{Associativitat.} $h\comp(g\comp f)=(h\comp g)\comp f$, per a $f\colon A\to B$, $g\colon B\to C$, $h\colon C\to D$.
\item \textbf{Identitats.} $\id_B\comp f=f$ i $f\comp\id_A=f$, per a tot $f\colon A\to B$.
\end{enumerate}
La secció~\ref{sec:diagrames} dona la lectura diagramàtica d'aquests axiomes.
\end{definicio}

Parlem aquí de \emph{col·leccions} i no de \emph{conjunts}, deliberadament: en alguns exemples els objectes, o els morfismes, són massa nombrosos per formar un conjunt. Així, tant $\Ob(\cat{C})$ com $\Mor(\cat{C})$ són en general col·leccions, no necessàriament conjunts. Precisarem aquesta qüestió a la secció~\ref{sec:mida}. Sovint és més còmode agrupar els morfismes segons el seu domini i codomini: quan la col·lecció de morfismes de $A$ a $B$ és un conjunt, se sol denotar
\[
\Hom_{\cat{C}}(A,B)
\]
o simplement $\Hom(A,B)$ quan la categoria és clara i l'anomenem \emph{homconjunt} de $A$ i $B$. Com que cada morfisme té un únic domini i un únic codomini, $\Mor(\cat{C})$ sempre es reparteix segons el parell $(A,B)$ de domini i codomini, en parts disjuntes dos a dos;\footnote{Aquesta disjunció és, de fet, equivalent a la unicitat del domini i el codomini. A l'axiomàtica relacional de l'apèndix~\ref{cap:axiomatica-relacional} en fa el paper d'un axioma.} el que no està garantit és que cada part sigui un \emph{conjunt}. Quan ho és per a tots els parells ---propietat que rep el nom de \emph{petitesa local} (secció~\ref{sec:mida})---, aquestes parts són els homconjunts $\Hom(A,B)$ i $\Mor(\cat{C})$ n'és la unió. Al llarg del llibre totes les categories que fem servir són localment petites, de manera que aquesta subtilesa no ens afectarà.

\begin{observacio}
Fem algunes precisions sobre la definició.
\begin{enumerate}
\item Els morfismes també s'anomenen \textbf{fletxes}, i sovint farem servir les dues paraules indistintament.
\item L'axioma d'associativitat fa que puguem escriure $h\comp g\comp f$ sense parèntesis, sabent que el resultat no depèn de la manera d'agrupar-los.
\item Els axiomes d'identitat diuen que $\id_A$ es comporta com un neutre per la composició. Veurem tot seguit que aquesta propietat determina $\id_A$ de manera única.
\item Aquesta definició axiomàtica pren com a termes no definits els objectes i els morfismes. Una categoria és qualsevol estructura que contingui aquestes dues menes de dades i que satisfaci les condicions i propietats de la definició.\footnotemark
\end{enumerate}
\end{observacio}
\footnotetext{De fet, es pot prescindir d'una de les dues menes de dades: com que cada objecte queda determinat per la seva identitat, una categoria petita es pot descriure en un llenguatge amb una sola mena d'individus ---les fletxes--- i tres relacions primitives de domini, codomini i composició, amb els objectes reconstruïts com les fletxes identitat. L'apèndix~\ref{cap:axiomatica-relacional} desenvolupa aquesta presentació relacional i en demostra l'equivalència amb la definició d'aquí.}

\begin{proposicio}
La fletxa identitat d'un objecte és única.
\end{proposicio}

\begin{prova}
Suposem que $e_A\colon A\to A$ i $e'_A\colon A\to A$ són dues fletxes identitat de $A$. Com que $e_A$ és una identitat,
\[
e_A\comp e'_A=e'_A.
\]
D'altra banda, com que $e'_A$ és una identitat,
\[
e_A\comp e'_A=e_A.
\]
Per tant $e_A=e'_A$.
\end{prova}

\section{Fletxes, camins i diagrames commutatius}\label{sec:diagrames}\index{diagrama commutatiu}

La definició anterior es pot llegir algebraicament, mitjançant equacions, però també gràficament. Aquesta segona lectura serà una part essencial del llenguatge categòric.

Si $f\colon A\to B$ i $g\colon B\to C$ són componibles, la composició es pot representar amb el triangle
\[
\begin{tikzpicture}[baseline=(m.center)]
  \matrix (m) [matrix of math nodes, column sep=3.4em, row sep=2.7em] {
    A & B \\
      & C \\
  };
  \draw[-{Stealth[scale=1.1]}] (m-1-1) -- node[above] {$f$} (m-1-2);
  \draw[-{Stealth[scale=1.1]}] (m-1-2) -- node[right] {$g$} (m-2-2);
  \draw[-{Stealth[scale=1.1]}] (m-1-1) -- node[below left] {$g\comp f$} (m-2-2);
\end{tikzpicture}
\]
La fletxa directa de $A$ a $C$ és precisament el resultat de seguir el camí $A\to B\to C$.

Més generalment, un diagrama pot contenir diferents camins dirigits amb el mateix origen i el mateix destí. Per exemple,
\[
\begin{tikzpicture}[baseline=(m.center)]
  \matrix (m) [matrix of math nodes, row sep=2.7em, column sep=3.5em] {
    A & B \\
    C & D \\
  };
  \draw[-{Stealth[scale=1.1]}] (m-1-1) -- node[above] {$f$} (m-1-2);
  \draw[-{Stealth[scale=1.1]}] (m-1-1) -- node[left] {$u$} (m-2-1);
  \draw[-{Stealth[scale=1.1]}] (m-1-2) -- node[right] {$v$} (m-2-2);
  \draw[-{Stealth[scale=1.1]}] (m-2-1) -- node[below] {$g$} (m-2-2);
\end{tikzpicture}
\]
té dos camins de $A$ a $D$: primer $f$ i després $v$, o primer $u$ i després $g$.

\begin{definicio}
Un diagrama com l'anterior \textbf{commuta} si els dos camins determinen el mateix morfisme, és a dir, si
\[
v\comp f=g\comp u.
\]
En general, direm que un diagrama commuta quan tots els camins dirigits amb el mateix origen i el mateix destí que cal comparar donen la mateixa composició.
\end{definicio}

\begin{observacio}
Un dibuix de fletxes no és automàticament un diagrama commutatiu. El dibuix especifica objectes i morfismes; afirmar que \emph{commuta} afegeix una o més igualtats entre composicions de fletxes. Per això convé aprendre a passar en tots dos sentits:
\[
\text{diagrama}\quad\longleftrightarrow\quad\text{equacions entre fletxes i composicions}.
\]
\end{observacio}

L'associativitat es pot llegir com un diagrama commutatiu. En una cadena
\[
A\xrightarrow{f}B\xrightarrow{g}C\xrightarrow{h}D
\]
podem agrupar les composicions de dues maneres, $(h\comp g)\comp f$ o $h\comp(g\comp f)$. El diagrama següent ho mostra com un paral·lelogram de vèrtexs $A,B,D,C$ i costats $f$, $h\comp g$, $h$ i $g\comp f$; la diagonal interior és $g$.
\[
\begin{tikzpicture}[>={Stealth[scale=1.1]}]
  \node (A) at (-3,2) {$A$};
  \node (B) at (0,2)  {$B$};
  \node (C) at (0,0)  {$C$};
  \node (D) at (3,0)  {$D$};
  \draw[->] (A) -- node[above] {$f$} (B);
  \draw[->] (B) -- node[right] {$h\comp g$} (D);
  \draw[->] (C) -- node[below] {$h$} (D);
  \draw[->] (A) -- node[below left=-2pt] {$g\comp f$} (C);
  \draw[->] (B) -- node[left] {$g$} (C);
\end{tikzpicture}
\]
Recórrer $A\to B\to C\to D$ agrupant primer $g\comp f$ i després component amb $h$, o agrupant primer $h\comp g$ i precomponent amb $f$, dona la mateixa fletxa $A\to D$. Això és exactament l'axioma d'associativitat:
\[
(h\comp g)\comp f=h\comp(g\comp f).
\]

Les identitats també tenen una lectura diagramàtica. Si $f\colon A\to B$, inserir $\id_A$ abans de $f$ o $\id_B$ després de $f$ no canvia el camí:
\[
\begin{tikzpicture}[baseline=(m.center)]
  \matrix (m) [matrix of math nodes, row sep=2.6em, column sep=2.6em] {
    A & A \\
      & B \\
  };
  \draw[-{Stealth[scale=1.1]}] (m-1-1) -- node[above] {\(\id_A\)} (m-1-2);
  \draw[-{Stealth[scale=1.1]}] (m-1-2) -- node[right] {\(f\)} (m-2-2);
  \draw[-{Stealth[scale=1.1]}] (m-1-1) -- node[below left] {\(f\)} (m-2-2);
\end{tikzpicture}
\qquad\qquad
\begin{tikzpicture}[baseline=(m.center)]
  \matrix (m) [matrix of math nodes, row sep=2.6em, column sep=2.6em] {
    A & B \\
      & B \\
  };
  \draw[-{Stealth[scale=1.1]}] (m-1-1) -- node[above] {\(f\)} (m-1-2);
  \draw[-{Stealth[scale=1.1]}] (m-1-2) -- node[right] {\(\id_B\)} (m-2-2);
  \draw[-{Stealth[scale=1.1]}] (m-1-1) -- node[below left] {\(f\)} (m-2-2);
\end{tikzpicture}
\]
Els dos triangles commutatius expressen exactament els axiomes d'identitat:
\[
f\comp\id_A=f,\qquad \id_B\comp f=f.
\]
A partir d'ara seguirem aquesta norma: quan una propietat categòrica tingui una expressió diagramàtica natural, en donarem alhora, al mateix lloc, l'equació i el diagrama. Llegir l'una i dibuixar l'altre, i a l'inrevés, és una destresa que el llibre demana constantment.

\section{Exemples de categories}

La força de la definició es veu de seguida en la varietat d'exemples que hi encaixen. En cadascun cal identificar els objectes, els morfismes, les identitats i la composició. Comencem pels exemples més petits, que malgrat la seva simplicitat apareixeran sovint com a \emph{categories índex} (la definició~\refalt{def:diagrama}{de diagrama del capítol~5}) per construir diagrames.

\subsection{Categories finites}\label{subsec:finites}\index{categoria!buida}\index{categoria!unitat}\index{categoria!finita}

Per descomptat, els objectes d'una categoria no cal que siguin conjunts. Considerem alguns exemples molt senzills que es poden descriure per complet. En cada cas identifiquem els objectes, els morfismes, la composició i les identitats; l'associativitat i les lleis d'identitat es comproven directament, ja que hi ha molt poques composicions a considerar.
\begin{enumerate}
\item La \emph{categoria buida} $\cat{0}$ té $\Ob(\cat{0})=\varnothing$ i cap morfisme. No hi ha composició ni identitats que definir, i els axiomes es compleixen trivialment perquè no hi ha res a comprovar; és la categoria \enquote{sense res}.
\item La \emph{categoria unitat} $\cat{1}$, també anomenada categoria \emph{terminal}, té $\Ob(\cat{1})=\{\ast\}$ i $\Hom(\ast,\ast)=\{\id_\ast\}$. L'única composició és $\id_\ast\comp\id_\ast=\id_\ast$, i els axiomes es redueixen a aquesta igualtat:
\[
\begin{tikzpicture}[baseline=(a.base)]
  \node (a) {$\ast$};
  \draw[-{Stealth[scale=1.1]}] (a) to[out=40,in=-40,looseness=8] node[right] {$\id_\ast$} (a);
\end{tikzpicture}
\]
\item La categoria $\cat{2}$ té $\Ob(\cat{2})=\{\ast,\star\}$ i una única fletxa no trivial $u\colon\ast\to\star$, a més de les dues identitats:
\[
\begin{aligned}
\Hom(\ast,\ast)&=\{\id_\ast\}, &
\Hom(\star,\star)&=\{\id_\star\},\\
\Hom(\ast,\star)&=\{u\}, &
\Hom(\star,\ast)&=\varnothing.
\end{aligned}
\]
Les úniques composicions són les que involucren identitats, $u\comp\id_\ast=u=\id_\star\comp u$, de manera que els axiomes es compleixen automàticament. A partir d'ara, com és costum, no dibuixem les identitats:
\[
\begin{tikzpicture}[baseline=(a.base)]
  \node (a) at (0,0) {$\ast$};
  \node (b) at (1.8,0) {$\star$};
  \draw[-{Stealth[scale=1.1]}] (a) -- node[above] {$u$} (b);
\end{tikzpicture}
\]
\item La categoria amb \textbf{dues fletxes paral·leles} té $\Ob=\{\ast,\star\}$ i dues fletxes no trivials $f,g\colon\ast\to\star$ amb el mateix domini i el mateix codomini:
\[
\Hom(\ast,\star)=\{f,g\},\qquad
\Hom(\star,\ast)=\varnothing,
\]
i als homconjunts $\Hom(\ast,\ast)$ i $\Hom(\star,\star)$ només hi ha les identitats. Com que cap parell d'aquestes fletxes no és componible, de nou les úniques composicions són amb identitats:
\[
\begin{tikzpicture}[baseline=(a.base)]
  \node (a) at (0,0) {$\ast$};
  \node (b) at (1.8,0) {$\star$};
  \draw[-{Stealth[scale=1.1]}] ([yshift=2.2pt]a.east) -- node[above] {$f$} ([yshift=2.2pt]b.west);
  \draw[-{Stealth[scale=1.1]}] ([yshift=-2.2pt]a.east) -- node[below] {$g$} ([yshift=-2.2pt]b.west);
\end{tikzpicture}
\]
Aquesta darrera reapareixerà com a categoria índex dels \emph{equalitzadors} (la definició~\refalt{def:equalitzador}{d'equalitzador del capítol~5}).
\item La categoria $\cat{3}$ té $\Ob(\cat{3})=\{\ast,\star,\bullet\}$ i tres fletxes no trivials, $f\colon\ast\to\star$, $g\colon\star\to\bullet$ i $h\colon\ast\to\bullet$:
\[
\Hom(\ast,\star)=\{f\},\qquad
\Hom(\star,\bullet)=\{g\},\qquad
\Hom(\ast,\bullet)=\{h\};
\]
els homconjunts d'un objecte a si mateix contenen només la identitat, i els homconjunts en sentit contrari són buits. La composició obliga a incloure la composició de les dues fletxes consecutives: l'única composició no trivial és $g\comp f=h$. Amb aquesta composició (i les identitats) els axiomes se satisfan:
\[
\begin{tikzpicture}[baseline=(m.center)]
  \matrix (m) [matrix of math nodes, row sep=2.6em, column sep=2.6em] {
    \ast & \star \\
         & \bullet \\
  };
  \draw[-{Stealth[scale=1.1]}] (m-1-1) -- node[above] {$f$} (m-1-2);
  \draw[-{Stealth[scale=1.1]}] (m-1-2) -- node[right] {$g$} (m-2-2);
  \draw[-{Stealth[scale=1.1]}] (m-1-1) -- node[below left] {$h=g\comp f$} (m-2-2);
\end{tikzpicture}
\]
\end{enumerate}

\begin{observacio}
Afegir fletxes no és gratuït: la composició obliga a incloure totes les composicions. Dues fletxes de sentits oposats $f\colon\ast\to\star$ i $g\colon\star\to\ast$
\[
\begin{tikzpicture}[baseline=(a.base)]
  \node (a) at (0,0) {$\ast$};
  \node (b) at (2.2,0) {$\star$};
  \draw[-{Stealth[scale=1.1]}] ([yshift=2.6pt]a.east) -- node[above] {$f$} ([yshift=2.6pt]b.west);
  \draw[-{Stealth[scale=1.1]}] ([yshift=-2.6pt]b.west) -- node[below] {$g$} ([yshift=-2.6pt]a.east);
\end{tikzpicture}
\]
generen $g\comp f$, $f\comp g$, $g\comp f\comp g$, i així indefinidament; la categoria resultant és infinita. Per obtenir-ne una de finita cal imposar equacions entre composicions ---per exemple $g\comp f=\id_\ast$ i $f\comp g=\id_\star$, cosa que fa de $f$ un isomorfisme (la definició~\ref{def:isomorfisme}). En canvi, dues fletxes \emph{paral·leles} $f,g\colon\ast\to\star$ no són componibles entre elles, i la categoria és finita sense afegir res.
\end{observacio}

\subsection{Categories de conjunts}\label{subsec:conjunts}

La categoria $\cat{Set}$ té:
\begin{itemize}
\item objectes: $\Ob(\cat{Set})$ és la col·lecció de tots els conjunts;
\item morfismes: $\Hom_{\cat{Set}}(A,B)$ és el conjunt de totes les aplicacions $A\to B$;
\item composició: la composició habitual d'aplicacions, $(g\comp f)(x)=g(f(x))$;
\item identitats: l'aplicació identitat, $\id_A(x)=x$.
\end{itemize}
L'associativitat i les lleis d'identitat són propietats ben conegudes de les aplicacions, de manera que $\cat{Set}$ és una categoria. Serà una categoria de referència: quan un concepte categòric nou ens desconcerti, mirar què vol dir a $\cat{Set}$ sol ser un bon primer pas.

Si ens restringim als conjunts finits i a les aplicacions entre ells, obtenim la categoria $\cat{FinSet}$: $\Ob(\cat{FinSet})$ és la col·lecció dels conjunts finits, $\Hom_{\cat{FinSet}}(A,B)=\Hom_{\cat{Set}}(A,B)$, i la composició i les identitats són les mateixes que a $\cat{Set}$. Com que la composició d'aplicacions entre conjunts finits torna a ser una aplicació entre conjunts finits, els axiomes se satisfan igual que a $\cat{Set}$.

Tot conjunt $S$ es pot veure també com una categoria, que continuem escrivint $S$:\index{categoria!discreta} $\Ob(S)=S$ i els únics morfismes són les identitats, una per objecte,
\[
\Hom_S(x,y)=\begin{cases}\{\id_x\} & \text{si } x=y,\\[2pt] \varnothing & \text{altrament.}\end{cases}
\]
La composició és forçada ($\id_x\comp\id_x=\id_x$) i els axiomes es compleixen trivialment. Una categoria en què els únics morfismes són les identitats s'anomena \textbf{categoria discreta}. En una categoria discreta no hi ha cap informació més enllà de la col·lecció d'objectes.

\begin{observacio}
Una categoria no queda determinada pels seus objectes: cal dir també quines són les fletxes. Amb només dos objectes, per exemple, hem trobat ja tres categories ben diferents:
\begin{itemize}
\item la \textbf{categoria discreta} de dos punts, només amb les identitats i cap fletxa no trivial;
\item la categoria $\cat{2}$, amb una única fletxa no trivial d'un objecte a l'altre;
\item la categoria amb \textbf{dues fletxes paral·leles}, amb dues fletxes no trivials entre els mateixos objectes.
\end{itemize}
Totes tres tenen els mateixos dos objectes, però són categories diferents perquè difereixen en les fletxes. En particular, $\cat{2}$ \emph{no} és la categoria discreta de dos punts: la presència de la fletxa no trivial les distingeix.
\end{observacio}

\subsection{Estructures elementals com a categories}\label{subsec:elementals}

Les estructures matemàtiques més simples sovint són categories. Per exemple, un monoide no s'assembla a una categoria d'un sol objecte: és exactament una categoria d'un sol objecte. Recollim aquesta idea en una llista d'identificacions, i tot seguit les justifiquem:
\begin{itemize}
\item un conjunt és una categoria discreta (ja vist a la subsecció anterior);
\item un monoide és una categoria d'un sol objecte;
\item un grup és una categoria d'un sol objecte en què tot morfisme és un isomorfisme;
\item un preordre és una categoria prima;
\item un ordre parcial és una categoria prima en què els únics isomorfismes són les identitats;
\item un sistema deductiu és una categoria prima.
\end{itemize}
En tots els casos, els objectes i els morfismes deixen de ser conjunts i aplicacions: passen a ser els elements i les relacions internes de l'estructura.

\paragraph{Preordres.}\index{preordre}
Un \emph{preordre} és un parell $(P,\le)$ on $P$ és un conjunt i $\le$ una relació binària sobre $P$ reflexiva ($x\le x$) i transitiva (si $x\le y$ i $y\le z$, aleshores $x\le z$). Per exemple, $(\mathbb{R},\le)$, o els subconjunts d'un conjunt ordenats per la inclusió $\subseteq$.

Tot preordre $(P,\le)$ \textbf{és} una categoria, que continuem escrivint $P$:
\begin{itemize}
\item objectes: $\Ob(P)=P$, els elements del preordre;
\item morfismes: hi ha exactament una fletxa $x\to y$ quan $x\le y$,
\[
\Hom_P(x,y)=\begin{cases}\{(x,y)\} & \text{si } x\le y,\\[2pt] \varnothing & \text{altrament;}\end{cases}
\]
\item composició: $(y,z)\comp(x,y)=(x,z)$, que existeix per transitivitat;
\item identitats: $\id_x=(x,x)$, que existeix per reflexivitat.
\end{itemize}
\[
\begin{tikzpicture}[baseline=(m.center)]
  \matrix (m) [matrix of math nodes, row sep=2.6em, column sep=2.6em] {
    x & y \\
      & z \\
  };
  \draw[-{Stealth[scale=1.1]}] (m-1-1) -- node[above] {$(x,y)$} (m-1-2);
  \draw[-{Stealth[scale=1.1]}] (m-1-2) -- node[right] {$(y,z)$} (m-2-2);
  \draw[-{Stealth[scale=1.1]}] (m-1-1) -- node[below left] {$(x,z)$} (m-2-2);
\end{tikzpicture}
\]
Com que cada homconjunt té com a màxim un element, l'associativitat i les lleis d'identitat se satisfan automàticament. Una categoria en què $|\Hom(x,y)|\le 1$ per a tot parell d'objectes s'anomena \textbf{prima}\index{categoria!prima} (en anglès, \emph{thin category}).

Si el preordre és a més antisimètric ($x\le y$ i $y\le x$ impliquen $x=y$), és un \emph{ordre parcial}. Categòricament, l'antisimetria diu que els únics isomorfismes (la definició~\ref{def:isomorfisme}) són les identitats: si $x\to y$ i $y\to x$ existeixen totes dues, aleshores $x=y$.

\paragraph{Monoides.}\index{monoide}
Un \emph{monoide} és una terna $(M,\ast,e)$ on $\ast$ és una operació associativa sobre $M$ amb element neutre $e$. Per exemple, $(\mathbb{N},+,0)$, o les paraules sobre un alfabet amb la concatenació i la paraula buida.

Tot monoide \textbf{és} una categoria amb un sol objecte $\ast$, que escrivim $\cat{B}M$.\footnote{La notació $\cat{B}M$ prové de la topologia, on $BM$ designa l'\emph{espai classificador} del monoide (o grup) $M$; la categoria $\cat{B}M$ n'és el model categòric.} Concretament:
\begin{itemize}
\item objectes: $\Ob(\cat{B}M)=\{\ast\}$;
\item morfismes: l'únic homconjunt és tot el monoide, $\Hom_{\cat{B}M}(\ast,\ast)=M$;
\item composició: l'operació del monoide, $g\comp f=g\ast f$, on l'element de la dreta actua primer;
\item identitat: $\id_\ast=e$.
\end{itemize}
Els axiomes de categoria són exactament els de monoide. Els morfismes, aquí, no són aplicacions: són els elements de $M$.

\paragraph{Grups.}\index{grup}
Un \emph{grup} és una terna $(G,\ast,e)$ que és un monoide en què, a més, tot element $a$ té un invers $a^{-1}$ amb $a\ast a^{-1}=e=a^{-1}\ast a$. Per exemple, $(\mathbb{Z},+,0)$, o les permutacions d'un conjunt amb la composició.

Un grup \textbf{és} una categoria d'un sol objecte en què tot morfisme és un isomorfisme. Com que un grup és en particular un monoide, la categoria és la de l'apartat anterior, que escrivim igualment $\cat{B}G$: $\Ob(\cat{B}G)=\{\ast\}$, $\Hom_{\cat{B}G}(\ast,\ast)=G$, $g\comp f=g\ast f$ i $\id_\ast=e$. En efecte, l'invers de l'element $a$ és exactament el morfisme invers de la fletxa corresponent. La teoria de grups és, doncs, la part invertible de la teoria de monoides. Cada element del grup és un bucle sobre l'únic objecte, i la identitat $e$ n'és el bucle destacat:
\[
\begin{tikzpicture}[baseline=(o.base)]
  \node (o) {$\ast$};
  % pètals genèrics (elements del grup)
  \foreach \ang in {30,150,210,270,330}{
    \draw[-{Stealth[scale=0.9]},gray] (o) to[out=\ang+14,in=\ang-14,looseness=16] (o);
  }
  % pètal destacat: la identitat e
  \draw[-{Stealth[scale=1.0]},thick] (o) to[out=104,in=76,looseness=16] (o);
  \node at (0,1.32) {$e$};
\end{tikzpicture}
\]

\paragraph{Sistemes deductius.}\index{deduibilitat@deduïbilitat}\index{Prov@$\cat{Prov}_T$}
Fixem un \emph{sistema deductiu} $T$ i considerem la seva relació de deduïbilitat $\vdash_T$ entre fórmules, reflexiva ($p\vdash_T p$) i transitiva (de $p\vdash_T q$ i $q\vdash_T r$ se segueix $p\vdash_T r$). Aquesta relació determina una categoria, que escrivim $\cat{Prov}_T$ per distingir-la del sistema $T$ del qual prové:
\begin{itemize}
\item objectes: $\Ob(\cat{Prov}_T)$ és el conjunt de les fórmules $p,q,r,\dots$ de $T$;
\item morfismes: hi ha exactament una fletxa $p\to q$ quan $p\vdash_T q$,
\[
\Hom_{\cat{Prov}_T}(p,q)=\begin{cases}\{(p,q)\} & \text{si } p\vdash_T q,\\[2pt] \varnothing & \text{altrament;}\end{cases}
\]
\item composició: $(q,r)\comp(p,q)=(p,r)$, que existeix per transitivitat;
\item identitats: $\id_p=(p,p)$, que existeix per reflexivitat.
\end{itemize}
\[
\begin{tikzpicture}[baseline=(m.center)]
  \matrix (m) [matrix of math nodes, row sep=2.6em, column sep=2.6em] {
    p & q \\
      & r \\
  };
  \draw[-{Stealth[scale=1.1]}] (m-1-1) -- node[above] {$p\vdash_T q$} (m-1-2);
  \draw[-{Stealth[scale=1.1]}] (m-1-2) -- node[right] {$q\vdash_T r$} (m-2-2);
  \draw[-{Stealth[scale=1.1]}] (m-1-1) -- node[below left] {$p\vdash_T r$} (m-2-2);
\end{tikzpicture}
\]
Com el preordre, $\cat{Prov}_T$ és una categoria prima. A diferència de l'ordre parcial, però, la deduïbilitat no acostuma a ser antisimètrica: dues fórmules diferents poden ser interdeduïbles ($p\vdash_T q$ i $q\vdash_T p$ amb $p\neq q$), com $p\land q$ i $q\land p$. Categòricament, això vol dir que hi ha isomorfismes no trivials: fórmules distintes però isomorfes. Aquesta categoria només registra \emph{si} $q$ es dedueix de $p$, no \emph{de quina manera} es dedueix; l'observació següent presenta la versió que distingeix les proves.

\begin{observacio}[Les proves mateixes com a fletxes]\index{prova!com a morfisme}
Hi ha una interpretació lògica més fina en què una prova concreta
\[
\pi\colon p\longrightarrow q
\]
és ella mateixa un morfisme, i dues proves $\pi,\rho\colon p\to q$ poden ser morfismes diferents. La composició representa aleshores l'encadenament de proves i la identitat, la prova trivial d'una fórmula a partir d'ella mateixa. Passem, doncs, d'una sola fletxa (registrar \emph{si} hi ha deducció) a moltes fletxes (registrar \emph{quina} prova).

Per obtenir una categoria de proves completament definida cal fixar el sistema deductiu i precisar també quan dues derivacions es consideren iguals ---per exemple, mitjançant una equivalència adequada entre proves. No necessitem encara aquesta teoria; de moment distingirem entre la categoria prima $\cat{Prov}_T$ de la \emph{deduïbilitat} i aquesta interpretació més rica de les \emph{proves com a fletxes}. Hi tornarem quan el llenguatge categòric ens permeti fer-ho amb més precisió.
\end{observacio}

\subsection{Categories de conjunts estructurats}

Ara considerem les categories els objectes de les quals són conjunts amb alguna estructura, i els morfismes, les aplicacions que preserven aquesta estructura. En tots els casos la composició és la d'aplicacions i les identitats són les aplicacions identitat, de manera que l'associativitat i les lleis d'identitat s'hereten de les de $\cat{Set}$; només cal comprovar, en cada cas, que la composició d'aplicacions que respecten l'estructura torna a respectar-la, i que la identitat també ho fa.

\paragraph{Grups.}
La categoria $\cat{Grp}$ té:
\begin{itemize}
\item objectes: $\Ob(\cat{Grp})$ és la col·lecció de tots els grups;
\item morfismes: $\Hom_{\cat{Grp}}(G,H)$ és el conjunt dels homomorfismes de grups $G\to H$;
\item composició: la composició d'aplicacions;
\item identitats: l'aplicació identitat $\id_{G}$.
\end{itemize}
Un \emph{homomorfisme de grups} $h\colon G\to H$ és una aplicació entre els conjunts subjacents que preserva l'operació: $h(a\ast b)=h(a)\ast h(b)$ per a tot $a,b$ (d'on se segueix que preserva també el neutre i els inversos). Si $h\colon G\to H$ i $k\colon H\to K$ són homomorfismes, la composició també ho és, ja que
\[
k(h(a\ast b))=k(h(a)\ast h(b))=k(h(a))\ast k(h(b));
\]
i l'aplicació identitat preserva trivialment l'operació. Com que la composició i les identitats són les d'aplicacions entre conjunts, l'associativitat i les lleis d'identitat s'hereten de $\cat{Set}$.

\paragraph{Espais vectorials.}\index{espai vectorial}
Fixem un cos $\mathbb{K}$. La categoria $\cat{Vect}_{\mathbb K}$ té:
\begin{itemize}
\item objectes: $\Ob(\cat{Vect}_{\mathbb K})$ és la col·lecció dels espais vectorials sobre $\mathbb{K}$;
\item morfismes: $\Hom_{\mathbb K}(V,W)$ és el conjunt de les aplicacions lineals $V\to W$;
\item composició: la composició d'aplicacions;
\item identitats: l'aplicació identitat $\id_V$.
\end{itemize}
Un \emph{espai vectorial} sobre $\mathbb{K}$ és un conjunt $V$ amb una suma commutativa i associativa $V\times V\to V$ ---amb element neutre i oposats--- i una multiplicació per escalars $\mathbb{K}\times V\to V$ compatible amb les operacions. Una \emph{aplicació lineal} $f\colon V\to W$ és una aplicació entre els conjunts subjacents que preserva l'estructura: $f(u+v)=f(u)+f(v)$ i $f(\lambda v)=\lambda f(v)$. Si $f\colon V\to W$ i $g\colon W\to U$ són lineals, la composició també ho és, ja que
\[
g(f(u+v))=g(f(u))+g(f(v))\qquad\text{i}\qquad g(f(\lambda v))=\lambda\,g(f(v));
\]
i la identitat és lineal. Per tant la composició i les identitats són les d'aplicacions entre conjunts, i l'associativitat i les lleis d'identitat s'hereten de $\cat{Set}$.

\paragraph{Espais topològics.}\index{espai topològic}
La categoria $\cat{Top}$ té:
\begin{itemize}
\item objectes: $\Ob(\cat{Top})$ és la col·lecció de tots els espais topològics;
\item morfismes: $\Hom_{\cat{Top}}(X,Y)$ és el conjunt de les aplicacions contínues $X\to Y$;
\item composició: la composició d'aplicacions;
\item identitats: l'aplicació identitat $\id_{X}$.
\end{itemize}
Un \emph{espai topològic} és un parell $(X,\mathcal{T}_X)$ on $X$ és un conjunt i $\mathcal{T}_X$ una família de subconjunts de $X$ (els \emph{oberts}) tal que:
\begin{itemize}
\item $\varnothing, X\in\mathcal{T}_X$;
\item si $U,V\in\mathcal{T}_X$, aleshores $U\cap V\in\mathcal{T}_X$;
\item si $\{U_i\}_{i\in I}$ és una família qualsevol en $\mathcal{T}_X$, aleshores $\bigcup_{i\in I}U_i\in\mathcal{T}_X$.
\end{itemize}
Una \emph{aplicació contínua} $f\colon(X,\mathcal{T}_X)\to(Y,\mathcal{T}_Y)$ és una aplicació $f\colon X\to Y$ que preserva l'estructura en el sentit següent: $f^{-1}(U)\in\mathcal{T}_X$ per a tot $U\in\mathcal{T}_Y$. Si $f\colon X\to Y$ i $g\colon Y\to Z$ són contínues, la composició també ho és, ja que $(g\comp f)^{-1}(U)=f^{-1}(g^{-1}(U))$ és obert quan $U$ ho és; i la identitat és contínua, perquè $\id^{-1}(U)=U$. Per tant, un cop més, la composició i les identitats són les d'aplicacions entre conjunts, i l'associativitat i les lleis d'identitat s'hereten de $\cat{Set}$.

\paragraph{Monoides.}
La categoria $\cat{Mon}$ té:
\begin{itemize}
\item objectes: $\Ob(\cat{Mon})$ és la col·lecció de tots els monoides;
\item morfismes: $\Hom_{\cat{Mon}}(M,N)$ és el conjunt dels homomorfismes de monoides $M\to N$;
\item composició: la composició d'aplicacions;
\item identitats: l'aplicació identitat $\id_{M}$.
\end{itemize}
Un \emph{homomorfisme de monoides} $h\colon M\to N$ és una aplicació entre els conjunts subjacents que preserva l'operació i el neutre: $h(a\ast b)=h(a)\ast h(b)$ i $h(e_M)=e_N$. Si $h\colon M\to N$ i $k\colon N\to P$ són homomorfismes, la composició també ho és, ja que
\[
k(h(a\ast b))=k(h(a)\ast h(b))=k(h(a))\ast k(h(b))\qquad\text{i}\qquad k(h(e_M))=k(e_N)=e_P;
\]
i la identitat preserva trivialment l'operació i el neutre. Com que la composició i les identitats són les d'aplicacions entre conjunts, l'associativitat i les lleis d'identitat s'hereten de $\cat{Set}$.

\paragraph{Ordres parcials.}
La categoria $\cat{Pos}$ té:
\begin{itemize}
\item objectes: $\Ob(\cat{Pos})$ és la col·lecció de tots els ordres parcials;
\item morfismes: $\Hom_{\cat{Pos}}(P,Q)$ és el conjunt de les aplicacions monòtones $P\to Q$;
\item composició: la composició d'aplicacions;
\item identitats: l'aplicació identitat $\id_{P}$.
\end{itemize}
Una \emph{aplicació monòtona} $f\colon P\to Q$ és una aplicació entre els conjunts subjacents que preserva l'ordre: si $x\le y$, aleshores $f(x)\le f(y)$. Si $f\colon P\to Q$ i $g\colon Q\to R$ són monòtones, la composició també ho és, ja que
\[
x\le y \ \Longrightarrow\ f(x)\le f(y) \ \Longrightarrow\ g(f(x))\le g(f(y));
\]
i la identitat preserva trivialment l'ordre. Un cop més, la composició i les identitats són les d'aplicacions entre conjunts, i les lleis s'hereten de $\cat{Set}$.

Convé no confondre les dues mirades: un monoide concret \emph{és} una categoria (d'un sol objecte) i un ordre parcial concret \emph{és} una categoria (prima), mentre que $\cat{Mon}$ i $\cat{Pos}$ tenen \emph{tots} els monoides i \emph{tots} els ordres parcials per objectes, respectivament.

\begin{observacio}[Categories com a contextos matemàtics]
Cadascuna d'aquestes categories és l'\emph{àmbit} on es desenvolupa una teoria matemàtica: $\cat{Top}$ és el context de la topologia general, $\cat{Grp}$ el de la teoria de grups, $\cat{Vect}_{\mathbb K}$ el de l'àlgebra lineal, i així successivament. Aquí, la categoria la formen tots els objectes d'una mena; a les estructures elementals, en canvi, un sol objecte estructurat ja és una categoria.
\end{observacio}

\subsection{Conjunts i relacions}\index{Rel@$\cat{Rel}$}

Com hem avançat a la introducció (secció~\ref{sec:mirar-fora}), els conjunts admeten una segona categoria, on els morfismes no són aplicacions sinó \emph{relacions}. La categoria $\cat{Rel}$ té:
\begin{itemize}
\item objectes: $\Ob(\cat{Rel})$ és la col·lecció de tots els conjunts;
\item morfismes: $\Hom_{\cat{Rel}}(X,Y)$ és el conjunt de les relacions $R\subseteq X\times Y$ (amb la precisió sobre els extrems de l'observació~\ref{obs:extrems-implicits});
\item composició: per a $R\subseteq X\times Y$ i $S\subseteq Y\times Z$,
\[
S\comp R=\{(x,z)\mid \exists y\in Y.\ (x,y)\in R \ \text{i}\ (y,z)\in S\}\subseteq X\times Z;
\]
\item identitats: la relació diagonal $\id_X=\{(x,x)\mid x\in X\}$.
\end{itemize}
Els axiomes es comproven directament a partir de la definició. Per a l'associativitat, un parell $(x,w)$ pertany a les dues maneres d'agrupar tres relacions encadenades exactament quan hi ha termes intermedis que connecten $x$ amb $w$ passant per les tres relacions; com que la condició no depèn de l'agrupament, les dues composicions coincideixen. Per a les identitats, compondre amb la diagonal no afegeix cap pas efectiu, perquè obliga el terme intermedi a coincidir amb l'extrem. La verificació completa és el primer exercici resolt del capítol~1 de la part de Pràctica (exercici~\ref{exr:pr-rel}). Compondre relacions equival a aplicar un quantificador existencial sobre el terme intermedi: aquest és el primer pont amb la lògica que assenyalàvem a la introducció. A diferència del que passa a $\cat{Set}$, un morfisme $X\to Y$ de $\cat{Rel}$ no ha d'assignar necessàriament a cada element de $X$ un únic element de $Y$: pot relacionar-lo amb diversos elements de $Y$, o amb cap.

\begin{observacio}[Extrems implícits]\label{obs:extrems-implicits}
En rigor, un morfisme $X\to Y$ de $\cat{Rel}$ no és una relació $R$ tota sola, sinó la relació \emph{junt amb} els seus extrems: la terna $(X,Y,R)$ amb $R\subseteq X\times Y$. Aquesta precisió és necessària perquè els homconjunts siguin disjunts, tal com exigeix la definició de categoria. Per exemple, la relació buida $R=\varnothing$ està continguda en $X\times Y$ per a qualsevol parell $(X,Y)$; si no registréssim els extrems, un mateix objecte $\varnothing$ seria alhora morfisme entre infinites parelles d'objectes, amb domini i codomini indeterminats. Registrant $(X,Y,R)$, cada morfisme té un únic domini $X$ i un únic codomini $Y$. Escriurem simplement $R$ quan els extrems se sobreentenguin; és la mateixa convenció d'\emph{extrems implícits} que ja hem fet servir en codificar la fletxa d'un preordre o de $\cat{Prov}_T$ com el parell $(x,y)$, i que retrobarem tot seguit a les categories sobre i sota un objecte. No canvia cap de les construccions: només n'explicita la compatibilitat amb la definició inicial de categoria.
\end{observacio}

\subsection{Categories de grafs}\label{subsec:grafs}

\begin{definicio}\index{graf dirigit}
Un \textbf{graf dirigit} $G$ consta d'un conjunt $V_G$ de \textbf{vèrtexs}, un conjunt $E_G$ d'\textbf{arestes} i dues aplicacions
\[
s_G,t_G\colon E_G\longrightarrow V_G,
\]
que assignen a cada aresta el seu \textbf{origen} i el seu \textbf{destí}, respectivament.
\end{definicio}

\begin{definicio}\index{morfisme de grafs}
Un \textbf{morfisme de grafs} $F\colon G\to H$ és una parella d'aplicacions
\[
F_V\colon V_G\to V_H,
\qquad
F_E\colon E_G\to E_H,
\]
que preserven origen i destí:
\[
F_V\comp s_G=s_H\comp F_E,
\qquad
F_V\comp t_G=t_H\comp F_E.
\]
Equivalentment, han de commutar els dos quadrats següents:

\begin{center}
\begin{tikzpicture}[baseline=(m.center)]
  \matrix (m) [matrix of math nodes, row sep=2.5em, column sep=3.2em] {
    E_G & V_G \\
    E_H & V_H \\
  };
  \draw[-{Stealth[scale=1.0]}] (m-1-1) -- node[above] {$s_G$} (m-1-2);
  \draw[-{Stealth[scale=1.0]}] (m-1-1) -- node[left] {$F_E$} (m-2-1);
  \draw[-{Stealth[scale=1.0]}] (m-1-2) -- node[right] {$F_V$} (m-2-2);
  \draw[-{Stealth[scale=1.0]}] (m-2-1) -- node[below] {$s_H$} (m-2-2);
\end{tikzpicture}
\hspace{2.2em}
\begin{tikzpicture}[baseline=(n.center)]
  \matrix (n) [matrix of math nodes, row sep=2.5em, column sep=3.2em] {
    E_G & V_G \\
    E_H & V_H \\
  };
  \draw[-{Stealth[scale=1.0]}] (n-1-1) -- node[above] {$t_G$} (n-1-2);
  \draw[-{Stealth[scale=1.0]}] (n-1-1) -- node[left] {$F_E$} (n-2-1);
  \draw[-{Stealth[scale=1.0]}] (n-1-2) -- node[right] {$F_V$} (n-2-2);
  \draw[-{Stealth[scale=1.0]}] (n-2-1) -- node[below] {$t_H$} (n-2-2);
\end{tikzpicture}
\end{center}
\end{definicio}

Els grafs dirigits i els morfismes de grafs formen una categoria, que denotarem $\cat{Graph}$:\index{Graph@$\cat{Graph}$}
\begin{itemize}
\item objectes: $\Ob(\cat{Graph})$ és la col·lecció de tots els grafs dirigits;
\item morfismes: $\Hom_{\cat{Graph}}(G,H)$ és el conjunt dels morfismes de grafs $G\to H$;
\item composició: component a component; per a $F\colon G\to H$ i $F'\colon H\to K$,
\[
F'\comp F=(F'_V\comp F_V,\ F'_E\comp F_E);
\]
\item identitats: $\id_G=(\id_{V_G},\id_{E_G})$.
\end{itemize}
La composició torna a ser un morfisme de grafs, perquè les equacions es conserven:
\[
(F'_V\comp F_V)\comp s_G=F'_V\comp s_H\comp F_E=s_K\comp(F'_E\comp F_E),
\]
i anàlogament per a $t$. L'associativitat i les lleis d'identitat es compleixen component a component, perquè són les de $\cat{Set}$.

Aquest exemple és especialment rellevant perquè els grafs són l'esquelet combinatori dels diagrames: tenen vèrtexs i arestes dirigides, però encara no incorporen per si sols identitats, composicions ni axiomes categòrics.

Per fixar idees, vet aquí un graf dirigit concret amb tres vèrtexs $V_G=\{a,b,c\}$ i tres arestes:
\[
\begin{tikzpicture}[baseline=(m.center)]
  \matrix (m) [matrix of math nodes, row sep=2.6em, column sep=2.6em] {
    a & b \\
      & c \\
  };
  \draw[-{Stealth[scale=1.1]}] (m-1-1) -- node[above] {$e_1$} (m-1-2);
  \draw[-{Stealth[scale=1.1]}] (m-1-2) -- node[right] {$e_2$} (m-2-2);
  \draw[-{Stealth[scale=1.1]}] (m-1-1) -- node[below left] {$e_3$} (m-2-2);
\end{tikzpicture}
\]
Aquí $e_1$ va de $a$ a $b$, $e_2$ de $b$ a $c$ i $e_3$ de $a$ a $c$. A diferència d'una categoria, un graf no demana que hi hagi cap aresta que \enquote{representi la composició} d'$e_1$ i $e_2$: $e_3$ hi és perquè l'hem posada, no perquè cap axioma l'obligui. Aquesta és exactament la diferència entre un graf i una categoria: la categoria afegeix identitats, composició i els axiomes que les regeixen. Per a un desenvolupament complet dels grafs, amb nombrosos exemples i els homomorfismes de grafs, vegeu l'apèndix~\ref{cap:grafs}.

\subsection{Resum: objectes i morfismes dels exemples}

La taula següent recull, per a cada exemple d'aquesta secció, quines dades constitueixen els objectes i quines els morfismes.

\begin{center}
\footnotesize
\setlength{\tabcolsep}{6pt}
\begin{tabularx}{\linewidth}{@{}>{\raggedright\arraybackslash}X>{\raggedright\arraybackslash}X>{\raggedright\arraybackslash}X@{}}
\toprule
Categoria & Objectes & Morfismes \\
\midrule
$\cat{0},\cat{1},\cat{2},\cat{3}$ & marcadors abstractes & fletxes formals \\
$\cat{Set}$, $\cat{FinSet}$ & conjunts & aplicacions \\
categoria discreta & marcadors abstractes & només identitats \\
$\cat{Grp}$, $\cat{Vect}_{\mathbb K}$, $\cat{Top}$, $\cat{Mon}$, $\cat{Pos}$ & conjunts amb estructura & aplicacions que preserven l'estructura \\
$\cat{B}M$ (un monoide) & un únic objecte $\ast$ & elements de $M$ \\
un preordre $P$; $\cat{Prov}_T$ & elements de $P$; fórmules de $T$ & $x\le y$; $p\vdash_T q$ \\
$\cat{Rel}$ & conjunts & relacions \\
$\cat{Graph}$ & grafs $(V,E,s,t)$ & parelles compatibles d'aplicacions \\
\bottomrule
\end{tabularx}
\end{center}

Aquesta taula fa visible que la definició de categoria no pressuposa res sobre la naturalesa dels objectes ni dels morfismes: no cal que els objectes siguin conjunts, ni que els morfismes siguin aplicacions. En una categoria discreta, per exemple, les identitats són morfismes \emph{formals}: presentar-les com a aplicacions pressuposaria una realització concreta que no forma part de la definició. I a $\cat{B}M$ els morfismes són els elements de $M$, que en general tampoc no són aplicacions. La categoria $\cat{Rel}$ ho remarca des de l'altre costat: els seus objectes són conjunts, però els seus morfismes ---les relacions--- no són funcions.

Convé no confondre aquesta descripció \emph{estructural} de les dades amb la seva eventual \emph{codificació} conjuntista. En el marc de fonaments que adoptem, qualsevol d'aquestes dades ---inclosos els marcadors abstractes--- es pot representar per conjunts; però aquesta codificació és una tria externa, no una part de la definició de categoria. Per això la taula descriu què \emph{són} les dades i no si \enquote{són conjunts} o \enquote{són funcions}: aquesta darrera pregunta barreja la naturalesa estructural amb una representació particular.

\section{Isomorfismes}\index{isomorfisme}

Entre tots els morfismes d'una categoria, n'hi ha uns d'especialment importants: els que tenen invers. Capturen, en el llenguatge abstracte de les fletxes, la idea que dos objectes tenen la mateixa estructura des del punt de vista de la categoria.

\begin{definicio}\label{def:isomorfisme}
En una categoria $\cat{C}$, un morfisme $f\colon A\to B$ és un \textbf{isomorfisme} si existeix un morfisme $g\colon B\to A$ tal que
\[
g\comp f=\id_A
\qquad\text{i}\qquad
f\comp g=\id_B.
\]
El morfisme $g$ s'anomena \textbf{invers} de $f$. Si existeix un isomorfisme entre $A$ i $B$, diem que $A$ i $B$ són \textbf{isomorfs}, i escrivim $A\cong B$.
\end{definicio}

La situació es caracteritza per dues equacions, $g\comp f=\id_A$ i $f\comp g=\id_B$, que es llegeixen com dos triangles commutatius:
\[
\begin{tikzpicture}[baseline=(m.center)]
  \matrix (m) [matrix of math nodes, row sep=2.6em, column sep=2.6em] {
    A & B \\
    A & \\
  };
  \draw[-{Stealth[scale=1.1]}] (m-1-1) -- node[above] {$f$} (m-1-2);
  \draw[-{Stealth[scale=1.1]}] (m-1-2) -- node[right] {$g$} (m-2-1);
  \draw[-{Stealth[scale=1.1]}] (m-1-1) -- node[left] {$\id_A$} (m-2-1);
\end{tikzpicture}
\qquad\qquad
\begin{tikzpicture}[baseline=(m.center)]
  \matrix (m) [matrix of math nodes, row sep=2.6em, column sep=2.6em] {
    B & A \\
    B & \\
  };
  \draw[-{Stealth[scale=1.1]}] (m-1-1) -- node[above] {$g$} (m-1-2);
  \draw[-{Stealth[scale=1.1]}] (m-1-2) -- node[right] {$f$} (m-2-1);
  \draw[-{Stealth[scale=1.1]}] (m-1-1) -- node[left] {$\id_B$} (m-2-1);
\end{tikzpicture}
\]
Cada triangle diu que anar d'un objecte a l'altre amb $f$ o $g$ i tornar és el mateix que seguir la identitat.

\begin{proposicio}
L'invers d'un isomorfisme, si existeix, és únic.
\end{proposicio}

\begin{prova}
Suposem que $f\colon A\to B$ té dos inversos, $g$ i $g'$. Aleshores
\[
g=g\comp\id_B=g\comp(f\comp g')=(g\comp f)\comp g'=\id_A\comp g'=g'.
\]
Per tant $g=g'$. Com que l'invers és únic, el podem denotar sense ambigüitat $f^{-1}$.
\end{prova}

\begin{proposicio}
Sigui $\cat{C}$ una categoria.
\begin{enumerate}
\item Per a cada objecte $A$, la identitat $\id_A$ és un isomorfisme, amb $\id_A^{-1}=\id_A$.
\item Si $f$ és un isomorfisme, aleshores $f^{-1}$ també ho és, i $(f^{-1})^{-1}=f$.
\item Si $f\colon A\to B$ i $g\colon B\to C$ són isomorfismes, aleshores $g\comp f$ també ho és, i
\[
(g\comp f)^{-1}=f^{-1}\comp g^{-1}.
\]
\end{enumerate}
\end{proposicio}

\begin{prova}
(1) De $\id_A\comp\id_A=\id_A$ se segueix que $\id_A$ és el seu propi invers. (2) Les equacions $g\comp f=\id_A$ i $f\comp g=\id_B$ que fan de $g=f^{-1}$ l'invers de $f$ diuen també que $f$ és l'invers de $g$; per tant $f^{-1}$ és un isomorfisme i $(f^{-1})^{-1}=f$. (3) Comprovem que $f^{-1}\comp g^{-1}$ és l'invers de $g\comp f$:
\[
(g\comp f)\comp(f^{-1}\comp g^{-1})=g\comp(f\comp f^{-1})\comp g^{-1}=g\comp\id_B\comp g^{-1}=g\comp g^{-1}=\id_C,
\]
i anàlogament $(f^{-1}\comp g^{-1})\comp(g\comp f)=\id_A$.
\end{prova}

\begin{corollari}
Donada una categoria $\cat{C}$, definim sobre $\Ob(\cat{C})$ la relació $\cong_{\cat{C}}$ per: $A\cong_{\cat{C}}B$ si i només si existeix un isomorfisme $f\colon A\to B$. Aleshores $\cong_{\cat{C}}$ és una relació d'equivalència.
\end{corollari}

\begin{prova}
És reflexiva: $\id_A$ és un isomorfisme per (1), de manera que $A\cong_{\cat{C}}A$. És simètrica: si $f$ és un isomorfisme de $A$ a $B$, aleshores $f^{-1}$ ho és de $B$ a $A$ per (2), de manera que $B\cong_{\cat{C}}A$. És transitiva: si $f\colon A\to B$ i $g\colon B\to C$ són isomorfismes, aleshores $g\comp f$ també ho és per (3), de manera que $A\cong_{\cat{C}}C$.
\end{prova}

\begin{observacio}[Igualtat i isomorfisme]\label{obs:igualtat-isomorfisme}\index{isomorfisme!i igualtat}
Convé fixar des d'ara una convenció que el llibre manté sempre. Escrivim $A=B$ només quan $A$ i $B$ són \emph{el mateix} objecte, i $A\cong B$ quan hi ha un isomorfisme entre ells. Quan un isomorfisme és distingit ---no depèn de cap tria---, l'anomenem \textbf{canònic}, però continua sent un isomorfisme i no una igualtat (en veurem un primer cas en el monoide lliure, exemple~\ref{ex:monoide-lliure-cat}).

La relació $\cong_{\cat{C}}$ permet parlar d'una \emph{igualtat estructural relativa a $\cat{C}$}: dos objectes isomorfs a $\cat{C}$ no es poden distingir per cap propietat formulada només amb els morfismes, la composició i les identitats de $\cat{C}$, perquè l'isomorfisme la transporta d'un a l'altre (ho precisarem a l'observació~\ref{obs:invariancia}). Diem, doncs, que tenen la mateixa estructura \emph{des del punt de vista de $\cat{C}$}, i per això sovint la pregunta rellevant no és si $A$ i $B$ són iguals, sinó si són isomorfs a $\cat{C}$.

Aquesta igualtat no és absoluta: depèn dels morfismes que la categoria considera admissibles. Per exemple, el conjunt $\{0,1\}$ amb l'ordre discret i el mateix conjunt amb l'ordre $0<1$ no són isomorfs a $\cat{Pos}$, perquè cap bijecció monòtona no té inversa monòtona, però els seus conjunts subjacents sí que són isomorfs a $\cat{Set}$. Canviar de categoria canvia, doncs, quins objectes compten com a estructuralment iguals. Aquesta és una de les idees centrals de la teoria de categories: la forma d'un objecte sempre es mira dins d'una categoria determinada.
La jerarquia completa ---igualtat, isomorfisme, isomorfisme natural i equivalència---, amb un contraexemple per a cada pas, es tracta a l'apèndix~\refalt{cap:contraexemples}{de contraexemples del llibre complet}.
\end{observacio}

\begin{exemple}
A $\cat{Set}$, els isomorfismes són exactament les aplicacions bijectives. En efecte, una aplicació té inversa per la composició si i només si és injectiva i exhaustiva.
\end{exemple}

\begin{exemple}[Isomorfismes en les categories de conjunts estructurats]
A $\cat{Grp}$, $\cat{Vect}_{\mathbb K}$ i $\cat{Mon}$, els isomorfismes són les aplicacions bijectives que preserven l'estructura i la inversa de les quals també la preserva; recuperem així les nocions habituals d'isomorfisme de grups, d'espais vectorials i de monoides. A $\cat{Top}$, en canvi, un isomorfisme és un homeomorfisme\index{homeomorfisme}: una aplicació contínua i bijectiva amb inversa contínua. Aquí es veu que la bijectivitat no basta: hi ha aplicacions contínues i bijectives que no són isomorfismes, perquè la seva inversa no és contínua. La definició categòrica capta la noció correcta sense haver-la d'enunciar cas per cas.
\end{exemple}

\begin{exemple}
En un preordre, un isomorfisme entre $x$ i $y$ vol dir que hi ha fletxes $x\to y$ i $y\to x$, és a dir,
\[
x\le y
\qquad\text{i}\qquad
y\le x.
\]
En un ordre parcial això obliga $x=y$. En un preordre general, en canvi, poden existir objectes isomorfs diferents: són exactament els elements equivalents per la relació $x\sim y$ si $x\le y$ i $y\le x$.
\end{exemple}

\begin{exemple}[Isomorfisme i equivalència lògica]\label{ex:iso-equivalencia-logica}
A la categoria prima $\cat{Prov}_T$ d'un sistema deductiu $T$, dues fórmules $p$ i $q$ són isomorfes exactament quan
\[
p\vdash_T q
\qquad\text{i}\qquad
q\vdash_T p.
\]
Per tant, l'isomorfisme categòric es tradueix aquí en \textbf{interdeduïbilitat}. Aquest és un primer exemple de com una noció categòrica adquireix una lectura lògica natural.
\end{exemple}

\begin{observacio}
L'exemple dels preordres mostra per què la definició categòrica és la bona. Es podria pensar a definir un isomorfisme com un morfisme \enquote{bijectiu}, però això requeriria parlar dels elements dels objectes, i no sempre té sentit. A més, fins i tot quan en té, pot donar la noció equivocada: a $\cat{Pos}$ existeixen aplicacions monòtones bijectives que no són isomorfismes perquè la seva inversa no és monòtona. La definició mitjançant composició i identitats evita aquest parany.
\end{observacio}

\begin{observacio}[Propietats invariants per isomorfisme]\label{obs:invariancia}
Sigui $\cat{C}$ una categoria. Una propietat dels objectes de $\cat{C}$ és \textbf{invariant per isomorfisme} a $\cat{C}$ si, sempre que un objecte la té, tot objecte isomorf a ell a $\cat{C}$ també la té. Com la relació $\cong_{\cat{C}}$ de l'observació anterior, aquesta noció és sempre relativa a una categoria determinada. Una propietat definida només a partir de l'estructura categòrica de $\cat{C}$ ---els seus objectes, morfismes, composició i identitats, sense fer servir la igualtat entre objectes--- ha de ser invariant per isomorfisme a $\cat{C}$, perquè un isomorfisme transporta qualsevol afirmació d'aquesta mena d'un objecte a l'altre. Per això la invariància per isomorfisme és un criteri bàsic per reconèixer quan una propietat depèn de la forma de l'objecte en la categoria i no de la seva presentació concreta.

Aquesta noció dona un mètode útil. Per provar que dos objectes \emph{són} isomorfs, n'hi ha prou d'exhibir un isomorfisme entre ells. Provar que \emph{no} ho són sol ser més difícil, perquè caldria descartar tots els isomorfismes possibles; la via pràctica és trobar una propietat invariant per isomorfisme a la categoria en qüestió que un dels objectes tingui i l'altre no. Per exemple, a $\cat{Pos}$, siguin $A$ la cadena $a<b<c$ i $B$ l'ordre amb $x<y$ i $x<z$ però $y,z$ incomparables. Comprovar directament que no són isomorfs exigiria descartar totes les bijeccions monòtones amb inversa monòtona; en canvi, n'hi ha prou d'observar que \emph{estar totalment ordenat} és una propietat invariant per isomorfisme, que $A$ té i $B$ no. Per tant $A$ i $B$ no són isomorfs.

La implicació inversa no és automàtica: una propietat pot ser invariant per isomorfisme sense haver estat formulada internament en el llenguatge categòric. L'exemple anterior n'és una mostra: \emph{estar totalment ordenat} és invariant per isomorfisme a $\cat{Pos}$, però l'hem definit parlant dels elements de l'ordre, no de fletxes.
\end{observacio}

\begin{observacio}
La noció d'isomorfisme és autènticament categòrica: es formula només amb composició i identitats, de manera que té sentit a \emph{qualsevol} categoria, sense referència als elements dels objectes. Això la fa més abstracta i més general que les nocions d'isomorfisme pròpies de cada context ---bijeccions, isomorfismes de grups, homeomorfismes---, que en són casos particulars.

Aquesta generalitat es fa visible en un cas que ja coneixem. Havíem identificat els monoides amb les categories d'un sol objecte; ara podem identificar els grups amb exactament aquelles categories d'un sol objecte en què tota fletxa és un isomorfisme. Portant la mateixa idea a un nombre qualsevol d'objectes s'obté una noció d'importància considerable en les matemàtiques actuals: un \emph{grupoide} és una categoria en què tot morfisme és un isomorfisme.\index{grupoide}

El grupoide unifica dos casos que ja hem trobat. Un grup és un grupoide amb un sol objecte, com acabem de veure. I un preordre en què tot morfisme és un isomorfisme és, precisament, un en què $x\le y$ implica $y\le x$: una relació simètrica, és a dir, una \emph{relació d'equivalència}. Així, una relació d'equivalència no és res més que un grupoide prim: les classes d'isomorfisme dels seus objectes són les classes d'equivalència.
\end{observacio}

\begin{definicio}\index{endomorfisme}\index{automorfisme}
Un morfisme $f\colon A\to A$, amb domini i codomini iguals, s'anomena \textbf{endomorfisme} de $A$. Un endomorfisme que a més és un isomorfisme s'anomena \textbf{automorfisme} de $A$.
\end{definicio}

Els automorfismes d'un objecte $A$, amb la composició, formen un grup, el \textbf{grup d'automorfismes} de $A$, que escrivim $\Aut(A)$.\index{Aut@$\Aut$} Torna a aparèixer així la identificació d'abans: $\Aut(A)$ és el grup associat a la categoria d'un sol objecte que té els automorfismes de $A$ com a fletxes.

\section{Classes especials de morfismes}

Els isomorfismes són els morfismes invertibles: tenen un invers per tots dos costats. Sovint, però, un morfisme té invers només per un costat, o satisfà una propietat de cancel·lació més feble que la invertibilitat. Aquestes nocions ---monomorfismes, epimorfismes, seccions i retraccions--- són bàsiques. A $\cat{Set}$, els monomorfismes i els epimorfismes coincideixen amb les aplicacions injectives i exhaustives; com veurem, però, aquesta coincidència és pròpia de $\cat{Set}$ i d'algunes categories concretes, no una propietat de les categories en general.

Recordem que, a $\cat{Set}$, una aplicació $f$ és \emph{injectiva} si $f(a)=f(a')$ implica $a=a'$, i \emph{exhaustiva} si tot element del codomini té alguna antiimatge. Aquestes definicions parlen d'elements; les versions categòriques que segueixen les reformulen només amb fletxes.

\subsection{Monomorfismes i epimorfismes}\label{subsec:mono-epi}

\begin{definicio}\index{monomorfisme}\index{epimorfisme}
\label{def:mono-epi}
Un morfisme $f\colon A\to B$ és un \textbf{monomorfisme} (o \textbf{mono}), i escrivim $f\colon A\rightarrowtail B$, si es pot cancel·lar per l'esquerra: per a tota parella de morfismes $g,h\colon C\to A$,
\[
\begin{tikzpicture}[baseline=(m.center)]
  \matrix (m) [matrix of math nodes, column sep=3.4em] { C & A & B \\ };
  \draw[-{Stealth[scale=1.0]}] ([yshift=2.6pt]m-1-1.east) -- node[above] {$g$} ([yshift=2.6pt]m-1-2.west);
  \draw[-{Stealth[scale=1.0]}] ([yshift=-2.6pt]m-1-1.east) -- node[below] {$h$} ([yshift=-2.6pt]m-1-2.west);
  \draw[-{Stealth[scale=1.1]}] (m-1-2) -- node[above] {$f$} (m-1-3);
\end{tikzpicture}
\qquad
f\comp g=f\comp h\ \Longrightarrow\ g=h.
\]
Un morfisme $f\colon A\to B$ és un \textbf{epimorfisme} (o \textbf{epi}), i escrivim $f\colon A\twoheadrightarrow B$, si es pot cancel·lar per la dreta: per a tota parella $g,h\colon B\to C$,
\[
\begin{tikzpicture}[baseline=(m.center)]
  \matrix (m) [matrix of math nodes, column sep=3.4em] { A & B & C \\ };
  \draw[-{Stealth[scale=1.1]}] (m-1-1) -- node[above] {$f$} (m-1-2);
  \draw[-{Stealth[scale=1.0]}] ([yshift=2.6pt]m-1-2.east) -- node[above] {$g$} ([yshift=2.6pt]m-1-3.west);
  \draw[-{Stealth[scale=1.0]}] ([yshift=-2.6pt]m-1-2.east) -- node[below] {$h$} ([yshift=-2.6pt]m-1-3.west);
\end{tikzpicture}
\qquad
g\comp f=h\comp f\ \Longrightarrow\ g=h.
\]
\end{definicio}

\begin{observacio}[Com cal llegir aquests diagrames]\label{obs:llegir-mono-epi}
Els diagrames de la definició~\ref{def:mono-epi} \emph{no} afirmen que res commuti: són esquemes de prova, i s'han de llegir per a cada parella possible de fletxes $g,h$. Al diagrama del monomorfisme hi ha dos camins de $C$ a $B$, \enquote{primer $g$, després $f$} i \enquote{primer $h$, després $f$}, i hi ha també les dues fletxes paral·leles $g,h$ de $C$ a $A$. Recordem que un diagrama amb dues fletxes paral·leles commuta exactament quan són iguals. La definició diu, doncs:
\begin{quote}
si els dos camins de $C$ a $B$ coincideixen, aleshores les fletxes paral·leles ja coincideixen.
\end{quote}
Amb el vocabulari de la secció~\ref{sec:diagrames}: l'antecedent no és que el diagrama commuti, sinó només que commutin els camins que acaben a $B$; la conclusió és que el diagrama commuta del tot. Un monomorfisme és, doncs, una fletxa per a la qual \emph{commutar després de $f$ implica commutar}.
Llegit en sentit contrari: si $g\neq h$, aleshores $f\comp g\neq f\comp h$. Un monomorfisme \emph{no confon} mai dues fletxes diferents que arriben a $A$; després de compondre-les amb $f$, continuen sent diferents.

Cal anar amb compte amb una lectura massa ràpida. Que $f$ sigui mono no impedeix que hi hagi parelles $g\neq h\colon C\to A$; el diagrama amb dues fletxes paral·leles diferents existeix perfectament. El que exclou és que, sent diferents, es facin iguals en compondre-les amb $f$.

El diagrama de l'epimorfisme es llegeix de la mateixa manera, amb les fletxes de prova a l'altre costat: si $g\neq h\colon B\to C$, aleshores $g\comp f\neq h\comp f$. Dues fletxes que surten de $B$ i coincideixen després de $f$ ja eren iguals; $f$ arriba prou a $B$ perquè cap diferència entre $g$ i $h$ no se li escapi. A $\cat{Set}$, prenent com a fletxes de prova les aplicacions des d'un conjunt d'un sol element, que són els elements de $A$, la lectura del monomorfisme es converteix exactament en la injectivitat, com veurem tot seguit.
\end{observacio}

\begin{definicio}\index{bimorfisme}
Un morfisme que és alhora un monomorfisme i un epimorfisme s'anomena \textbf{bimorfisme}.
\end{definicio}

\begin{exemple}
A $\cat{Set}$, els monomorfismes són exactament les aplicacions \textbf{injectives} i els epimorfismes, les \textbf{exhaustives}. Vegem el cas dels monomorfismes: si $f$ és injectiva i $f\comp g=f\comp h$, aleshores per a cada $c$ tenim $f(g(c))=f(h(c))$, d'on $g(c)=h(c)$; recíprocament, si $f$ no és injectiva, hi ha $a\neq a'$ amb $f(a)=f(a')$, i les dues aplicacions constants $g,h\colon\{\ast\}\to A$ amb valors $a$ i $a'$ compleixen $f\comp g=f\comp h$ però $g\neq h$.

Per als epimorfismes, sigui $f\colon A\to B$. Si $f$ és exhaustiva i $g\comp f=h\comp f$, aleshores per a cada $b\in B$ existeix $a$ amb $f(a)=b$, i $g(b)=g(f(a))=h(f(a))=h(b)$, d'on $g=h$. Recíprocament, si $f$ no és exhaustiva, hi ha $b_0\in B$ sense antiimatge, és a dir, $b_0\notin f(A)$. Definim $g,h\colon B\to\{0,1\}$ a \emph{tot} $B$ per
\[
g(b)=0 \text{ per a tot } b\in B,
\qquad
h(b)=\begin{cases}1 & \text{si } b=b_0,\\ 0 & \text{si } b\neq b_0.\end{cases}
\]
Aleshores $g$ i $h$ coincideixen a tot $B$ tret de $b_0$, de manera que $g\neq h$. Però per a tot $a\in A$ tenim $f(a)\neq b_0$, ja que $b_0\notin f(A)$; per tant $g(f(a))=0=h(f(a))$, és a dir, $g\comp f=h\comp f$. Això mostra que $f$ no és un epimorfisme.

Convé notar que cap dels dos arguments no tria una \emph{secció} de $f$: per als monomorfismes es comparen dues seleccions d'elements des d'un conjunt unitari, i per als epimorfismes es construeixen explícitament dues aplicacions cap a $\{0,1\}$. En particular, la caracterització \enquote{epimorfisme $=$ exhaustiva} a $\cat{Set}$ no fa servir l'axioma de l'elecció, i no s'ha de confondre amb l'enunciat, ben diferent, que tota aplicació exhaustiva \emph{admet una secció} ---aquest sí equivalent a l'axioma de l'elecció (vegeu més avall).
\end{exemple}

\begin{proposicio}
Tot isomorfisme és un bimorfisme: si $f$ és un isomorfisme, aleshores és alhora un monomorfisme i un epimorfisme.
\end{proposicio}

\begin{prova}
Sigui $f\colon A\to B$ un isomorfisme amb invers $f^{-1}$. Si $f\comp g=f\comp h$, aleshores
\[
g=f^{-1}\comp(f\comp g)=f^{-1}\comp(f\comp h)=h,
\]
de manera que $f$ és mono. Si $g\comp f=h\comp f$, aleshores
\[
g=(g\comp f)\comp f^{-1}=(h\comp f)\comp f^{-1}=h,
\]
de manera que $f$ és epi.
\end{prova}

\begin{observacio}
A $\cat{Set}$, un morfisme és un isomorfisme si i només si és alhora mono i epi, és a dir, si i només si, com a aplicació, és injectiu i exhaustiu, o sigui \textbf{bijectiu}. Dit d'una altra manera, a $\cat{Set}$ els bimorfismes són exactament els isomorfismes. El recíproc de la proposició anterior, però, no és cert: en general un bimorfisme no és un isomorfisme. L'única fletxa no trivial de la categoria $\cat{2}$ és un bimorfisme ---per la simple raó que gairebé no hi ha morfismes amb què fer les cancel·lacions---, però no té invers, de manera que no és un isomorfisme. Un exemple més substancial es dona a $\cat{Top}$: una bijecció contínua és un bimorfisme, però no és un isomorfisme si la seva inversa no és contínua.
\end{observacio}

\begin{observacio}[Atenció: raonar amb elements]\label{obs:raonar-elements}\index{element!perill de raonar amb}
A $\cat{Set}$ podem raonar amb elements per dos motius: cada element $a\in A$ és una fletxa $\{\ast\}\to A$, i dues aplicacions $A\to B$ són iguals si coincideixen sobre tots els elements. En una categoria qualsevol, cap d'aquestes dues coses no està garantida.
\begin{itemize}
\item \emph{Pot ser que no hi hagi elements.} En un preordre o en un monoide vist com a categoria, una fletxa no envia res enlloc, i preguntar si és \enquote{injectiva} no té sentit. Mono i epi, en canvi, sempre tenen sentit, perquè es defineixen per composició (exercici~\ref{exr:pr-prima-mono-epi} de la Pràctica).
\item \emph{L'aplicació subjacent pot enganyar.} A $\cat{Top}$, una bijecció contínua és mono i epi sense ser un isomorfisme (observació anterior). A $\cat{Mon}$, la inclusió $\iota\colon(\mathbb N,+,0)\hookrightarrow(\mathbb Z,+,0)$ és un epimorfisme, tot i que no és exhaustiva. Si dos homomorfismes $g,h\colon\mathbb Z\to M$ coincideixen sobre $\mathbb N$, aleshores $g(-1)$ i $h(-1)$ són tots dos inversos de $g(1)=h(1)$, i en un monoide l'invers d'un element, si existeix, és únic. Per tant, $g(-1)=h(-1)$ i, com que tot enter és suma d'un natural i de còpies de $-1$, $g=h$.
\item \emph{Els \enquote{punts} poden no separar fletxes.} A $\cat{Grp}$, el grup trivial $1$ fa el paper de $\{\ast\}$, però de $1$ a qualsevol grup només hi ha un homomorfisme. Les fletxes $1\to G$ no distingeixen, doncs, cap parell d'homomorfismes diferents.
\end{itemize}
Per això, en una categoria general els arguments es fan amb fletxes i composicions, i reservem els adjectius \emph{injectiu}, \emph{exhaustiu} i \emph{bijectiu} per a les aplicacions, incloses les aplicacions subjacents dels morfismes d'una categoria concreta. La definició de monomorfisme és, de fet, una \enquote{injectivitat} respecte de \emph{totes} les fletxes de prova $C\to A$, i no només respecte dels elements.
\end{observacio}

\subsection{Seccions i retraccions}\label{subsec:seccions-retraccions}

\begin{definicio}\index{secció}\index{retracció}\index{invers!per un costat}
Considerem un morfisme $f\colon A\to B$. Diem que un morfisme $s\colon B\to A$ és una \textbf{secció} de $f$ quan es compleix $f\comp s=\id_B$. Això equival a dir que el diagrama següent és commutatiu:
\[
\begin{tikzpicture}[baseline=(m.center)]
  \matrix (m) [matrix of math nodes, row sep=2.6em, column sep=2.6em] {
    B & A \\
    B & \\
  };
  \draw[-{Stealth[scale=1.1]}] (m-1-1) -- node[above] {$s$} (m-1-2);
  \draw[-{Stealth[scale=1.1]}] (m-1-2) -- node[right] {$f$} (m-2-1);
  \draw[-{Stealth[scale=1.1]}] (m-1-1) -- node[left] {$\id_B$} (m-2-1);
\end{tikzpicture}
\]
També diem en aquest cas que $s$ és un \emph{invers per la dreta} de $f$.

Simètricament, diem que un morfisme $r\colon B\to A$ és una \textbf{retracció} de $f$ quan es compleix $r\comp f=\id_A$. Això equival a dir que el diagrama següent és commutatiu:
\[
\begin{tikzpicture}[baseline=(m.center)]
  \matrix (m) [matrix of math nodes, row sep=2.6em, column sep=2.6em] {
    A & B \\
    A & \\
  };
  \draw[-{Stealth[scale=1.1]}] (m-1-1) -- node[above] {$f$} (m-1-2);
  \draw[-{Stealth[scale=1.1]}] (m-1-2) -- node[right] {$r$} (m-2-1);
  \draw[-{Stealth[scale=1.1]}] (m-1-1) -- node[left] {$\id_A$} (m-2-1);
\end{tikzpicture}
\]
També diem en aquest cas que $r$ és un \emph{invers per l'esquerra} de $f$.
\end{definicio}

Una conseqüència òbvia d'aquestes dues definicions és que quan es compleix, per exemple, $f\comp s=\id_B$, se'n segueix alhora que $s$ és una secció de $f$ i que $f$ és una retracció de $s$; el mateix succeeix amb $r\comp f=\id_A$, on $r$ és una retracció de $f$ i $f$ és una secció de $r$.

Un morfisme pot tenir cap, una o diverses seccions, i igualment cap, una o diverses retraccions. Vegem-ne exemples a $\cat{Set}$:
\begin{itemize}
\item L'única aplicació $f\colon\emptyset\to\{*\}$ no té cap secció: no existeix cap aplicació $\{*\}\to\emptyset$.
\item La injecció $f\colon\{0\}\hookrightarrow\{0,1\}$ associada a la inclusió $\{0\}\subseteq\{0,1\}$ té una única retracció, l'única aplicació $\{0,1\}\to\{0\}$; en canvi no té cap secció.
\item L'aplicació $f\colon\{0,1\}\to\{*\}$ té dues seccions distintes, $*\mapsto 0$ i $*\mapsto 1$: cada tria d'una antiimatge en dona una.
\end{itemize}
Tenir secció o retracció és, doncs, una propietat especial de $f$, i quan en té, en general no és única.

\begin{proposicio}\label{prop:escindit-mono-epi}
Sigui $f\colon A\to B$. Si $f$ té alguna secció, aleshores $f$ és un epimorfisme. Si $f$ té alguna retracció, aleshores $f$ és un monomorfisme.
\end{proposicio}

\begin{prova}
Suposem que $f$ té una retracció $r$, amb $r\comp f=\id_A$, i siguin $g,h\colon C\to A$ amb $f\comp g=f\comp h$. Component per l'esquerra amb $r$,
\[
g=\id_A\comp g=(r\comp f)\comp g=r\comp(f\comp g)=r\comp(f\comp h)=(r\comp f)\comp h=h,
\]
de manera que $f$ és mono. El cas de la secció és anàleg: si $f$ té secció $s$, amb $f\comp s=\id_B$, i $g,h\colon B\to C$ amb $g\comp f=h\comp f$, aleshores component per la dreta amb $s$,
\[
g=g\comp\id_B=g\comp(f\comp s)=(g\comp f)\comp s=(h\comp f)\comp s=h\comp(f\comp s)=h\comp\id_B=h,
\]
de manera que $f$ és epi.
\end{prova}

\begin{definicio}\index{morfisme!escindit}\index{epimorfisme!escindit}\index{monomorfisme!escindit}
Diem que un morfisme $f$ queda \textbf{escindit} quan té una secció o una retracció. En concret, si $f$ té una secció, aleshores $f$ és un epimorfisme i se'n diu \textbf{epimorfisme escindit}; si $f$ té una retracció, aleshores $f$ és un monomorfisme i se'n diu \textbf{monomorfisme escindit}.
\end{definicio}

\begin{observacio}
El recíproc no val en general: hi ha monomorfismes que no tenen cap retracció (i per tant no són escindits) i epimorfismes que no tenen cap secció. En efecte, a $\cat{Set}$ tota aplicació injectiva amb domini no buit té alguna retracció i tota exhaustiva té alguna secció ---aquesta darrera afirmació és, de fet, equivalent a l'\emph{axioma de l'elecció}---, però en altres categories la situació és diferent.

Finalment, un morfisme $f$ és un \textbf{isomorfisme} si i només si té alhora alguna secció i alguna retracció. En aquest cas la secció i la retracció són forçosament \emph{úniques} i coincideixen: si $f\comp s=\id_B$ i $r\comp f=\id_A$, aleshores
\[
r=r\comp\id_B=r\comp(f\comp s)=(r\comp f)\comp s=\id_A\comp s=s.
\]
En canvi, la unicitat d’una secció o d’una retracció, considerada
separadament, no caracteritza els isomorfismes en una categoria general; per exemple, la injecció
\[
 i\colon\{0\}\hookrightarrow\{0,1\}
\]
té una única retracció $r\colon\{0,1\}\to\{0\}$, però no és un isomorfisme. No és una subtilesa avançada, sinó una contradicció amb un exemple de la mateixa subsecció.
\end{observacio}

\section{Construccions sobre categories}

A partir de categories donades se'n poden construir de noves. Aquestes construccions seran útils constantment, i algunes ---com la categoria oposada--- són fonamentals.

\subsection{La categoria oposada}\index{categoria!oposada}

\begin{definicio}
Donada una categoria $\cat{C}$, la seva \textbf{categoria oposada} $\cat{C}\op$ té:
\begin{itemize}
\item objectes: els mateixos que $\cat{C}$, $\Ob(\cat{C}\op)=\Ob(\cat{C})$;
\item morfismes: per a cada morfisme $f\colon A\to B$ de $\cat{C}$, un morfisme $f\op\colon B\to A$ de $\cat{C}\op$, de manera que
\[
\Hom_{\cat{C}\op}(B,A)=\{f\op\mid f\in\Hom_{\cat{C}}(A,B)\};
\]
\item composició: s'obté invertint l'ordre; per a $f\colon A\to B$ i $g\colon B\to C$ de $\cat{C}$,
\[
f\op\comp g\op=(g\comp f)\op\colon C\to A;
\]
\item identitats: la identitat de $A$ a $\cat{C}\op$ és $(\id_A)\op$.
\end{itemize}
És a dir, $\cat{C}\op$ s'obté invertint el sentit de totes les fletxes.
\end{definicio}

Visualment, la composició a $\cat{C}$ i la corresponent a $\cat{C}\op$ inverteixen l'ordre:
\[
\begin{array}{c@{\qquad\qquad}c}
\cat{C}: & \cat{C}\op: \\[0.6em]
\begin{tikzpicture}[baseline=(m.center)]
  \matrix (m) [matrix of math nodes, row sep=2.6em, column sep=2.6em] {
    A & B \\
      & C \\
  };
  \draw[-{Stealth[scale=1.1]}] (m-1-1) -- node[above] {$f$} (m-1-2);
  \draw[-{Stealth[scale=1.1]}] (m-1-2) -- node[right] {$g$} (m-2-2);
  \draw[-{Stealth[scale=1.1]}] (m-1-1) -- node[below left] {$g\comp f$} (m-2-2);
\end{tikzpicture}
&
\begin{tikzpicture}[baseline=(m.center)]
  \matrix (m) [matrix of math nodes, row sep=2.6em, column sep=2.6em] {
    A & B \\
      & C \\
  };
  \draw[-{Stealth[scale=1.1]}] (m-1-2) -- node[above] {$f\op$} (m-1-1);
  \draw[-{Stealth[scale=1.1]}] (m-2-2) -- node[right] {$g\op$} (m-1-2);
  \draw[-{Stealth[scale=1.1]}] (m-2-2) -- node[below left] {$f\op\comp g\op$} (m-1-1);
\end{tikzpicture}
\end{array}
\]
L'associativitat a $\cat{C}\op$ es dedueix de la de $\cat{C}$ invertint l'ordre, i les lleis d'identitat, de les de $\cat{C}$: per exemple, $f\op\comp(\id_B)\op=(\id_B\comp f)\op=f\op$. La categoria oposada és la base de la \textbf{dualitat}\index{dualitat}, que desenvolupem a la secció~\ref{sec:dualitat}: tota afirmació formulada purament en el llenguatge categòric dona, invertint les fletxes, una afirmació dual, i aquesta correspondència es pot fer del tot precisa.

\begin{observacio}
Aplicar dues vegades l'oposada retorna la categoria de partida:
\[
(\cat{C}\op)\op=\cat{C}.
\]
Invertir les fletxes i tornar-les a invertir les deixa com estaven.
\end{observacio}

Vegem com queda l'oposada en dos casos. Quan un preordre $(P,\le)$ es considera com una categoria ---amb un morfisme $x\to y$ quan $x\le y$---, la seva oposada té un morfisme $x\to y$ exactament quan $y\le x$. I quan un monoide $(M,\ast,e)$ es considera com una categoria, la seva oposada té els mateixos morfismes ---els elements de $M$---, però la composició de $a$ i $b$ és $b\ast a$ en lloc de $a\ast b$.

\subsection{La categoria producte}\label{subsec:categoria-producte}\index{categoria!producte}

\begin{definicio}
Donades dues categories $\cat{C}$ i $\cat{D}$, la \textbf{categoria producte} $\cat{C}\times\cat{D}$ té:
\begin{itemize}
\item objectes: $\Ob(\cat{C}\times\cat{D})$ és la col·lecció de les parelles $(A,X)$ amb $A\in\Ob(\cat{C})$ i $X\in\Ob(\cat{D})$;
\item morfismes: $\Hom_{\cat{C}\times\cat{D}}\bigl((A,X),(B,Y)\bigr)=\Hom_{\cat{C}}(A,B)\times\Hom_{\cat{D}}(X,Y)$, és a dir, les parelles $(f,g)$ amb $f\colon A\to B$ a $\cat{C}$ i $g\colon X\to Y$ a $\cat{D}$;
\item composició: component a component,
\[
(f',g')\comp(f,g)=(f'\comp f,\ g'\comp g);
\]
\item identitats: $\id_{(A,X)}=(\id_A,\id_X)$.
\end{itemize} Els axiomes de categoria per a $\cat{C}\times\cat{D}$ se segueixen dels de $\cat{C}$ i $\cat{D}$ aplicats a cada component.
\end{definicio}

\begin{exemple}
Prenem la categoria $\cat{2}$ de la subsecció~\ref{subsec:finites} ---dos objectes i una única fletxa no trivial---, i n'escrivim els objectes com $0$ i $1$, amb la fletxa no trivial $0\to 1$. La categoria producte $\cat{2}\times\cat{2}$ té quatre objectes, que es disposen de manera natural en un quadrat:
\[
\begin{tikzpicture}[baseline=(m.center)]
  \matrix (m) [matrix of math nodes, row sep=2.6em, column sep=3.4em] {
    (0,0) & (0,1) \\
    (1,0) & (1,1) \\
  };
  \draw[-{Stealth[scale=1.1]}] (m-1-1) -- (m-1-2);
  \draw[-{Stealth[scale=1.1]}] (m-1-1) -- (m-2-1);
  \draw[-{Stealth[scale=1.1]}] (m-1-2) -- (m-2-2);
  \draw[-{Stealth[scale=1.1]}] (m-2-1) -- (m-2-2);
  \draw[-{Stealth[scale=1.1]}] (m-1-1) -- (m-2-2);
\end{tikzpicture}
\]
Cada fletxa és una parella de fletxes de $\cat{2}$. Com que $\cat{2}$ és prima, també ho és $\cat{2}\times\cat{2}$: entre dos objectes hi ha com a màxim un morfisme. Per això els dos camins del quadrat de $(0,0)$ a $(1,1)$ tenen necessàriament la mateixa composició, que és la diagonal; el quadrat commuta.
\end{exemple}

\begin{observacio}[Dos usos de \enquote{producte}]\label{obs:dos-productes}\index{producte!d'objectes i de categories}
No s'han de confondre dues construccions que porten el mateix nom. La categoria producte $\cat{C}\times\cat{D}$ és una \emph{categoria} nova, construïda a partir de dues categories. En canvi, el producte $A\times B$ de dos objectes ---que la introducció fa servir com a primer exemple de propietat universal i que avançarem a la subsecció~\ref{subsec:producte-coproducte}--- és un \emph{objecte} d'una mateixa categoria, caracteritzat per una propietat universal. Les dues nocions estan relacionades: quan disposem de la categoria $\cat{Cat}$ de les categories petites (capítol~\ref{cap:functors}), si $\cat{C}$ i $\cat{D}$ són petites, la categoria $\cat{C}\times\cat{D}$ es podrà veure com el producte dels objectes $\cat{C}$ i $\cat{D}$ dins de $\cat{Cat}$. Però no són la mateixa construcció: l'una fabrica una categoria a partir de dues, i l'altra viu dins d'una sola categoria.
\end{observacio}

\subsection{Subcategories}\label{subsec:subcategories}\index{subcategoria}

Una altra manera d'obtenir noves categories a partir de categories antigues és restringint-ne els objectes o les fletxes.

\begin{definicio}
Una \textbf{subcategoria} $\cat{S}$ d'una categoria $\cat{C}$ ve donada per una subcol·lecció d'objectes $\Ob(\cat{S})\subseteq\Ob(\cat{C})$ i, per a cada parella $A,B\in\Ob(\cat{S})$, una subcol·lecció de morfismes $\Hom_{\cat{S}}(A,B)\subseteq\Hom_{\cat{C}}(A,B)$, de manera que:
\begin{itemize}
\item la identitat $\id_A$ pertany a $\Hom_{\cat{S}}(A,A)$ per a cada $A\in\Ob(\cat{S})$;
\item la composició de dos morfismes de $\cat{S}$ componibles pertany a $\cat{S}$.
\end{itemize}
Amb aquestes condicions, $\cat{S}$ és una categoria per dret propi, amb la composició i les identitats heretades de $\cat{C}$.
\end{definicio}

Dos casos particulars, en cert sentit oposats, són especialment útils. Es diu que $\cat{S}$ és una \textbf{subcategoria plena}\index{subcategoria!plena} quan $\Hom_{\cat{S}}(A,B)=\Hom_{\cat{C}}(A,B)$ per a tot $A,B\in\Ob(\cat{S})$; una subcategoria plena queda determinada, doncs, només per la seva col·lecció d'objectes. Simètricament, es diu que $\cat{S}$ és una \textbf{subcategoria ampla}\index{subcategoria!ampla} quan $\Ob(\cat{S})=\Ob(\cat{C})$, encara que només contingui alguns dels morfismes. En resum: una subcategoria plena restringeix els objectes i conserva tots els morfismes entre ells; una subcategoria ampla conserva tots els objectes i en restringeix els morfismes.

\begin{exemple}
$\cat{FinSet}$ és una subcategoria plena de $\cat{Set}$: pren com a objectes els conjunts finits i, entre dos d'ells, totes les aplicacions. En canvi, la subcategoria de $\cat{Set}$ que té tots els conjunts per objectes però només les aplicacions \emph{injectives} per morfismes és una subcategoria ampla que no és plena: conserva tots els objectes, però deixa fora les aplicacions que no són injectives.
\end{exemple}

\subsection{Categories sobre i sota un objecte}\index{categoria!sobre un objecte}\index{categoria!sota un objecte}

Fixat un objecte $A$ d'una categoria $\cat{C}$, se'n pot construir una categoria nova els objectes de la qual són els morfismes que apunten cap a $A$.

\begin{definicio}\label{def:slice}
Sigui $\cat{C}$ una categoria i $A$ un objecte. La \textbf{categoria de $\cat{C}$ sobre $A$} (en anglès, \emph{slice category}), que escrivim $\cat{C}/A$, té:
\begin{itemize}
\item objectes: $\Ob(\cat{C}/A)$ és la col·lecció dels morfismes $x\colon X\to A$ de $\cat{C}$ amb codomini $A$;
\item morfismes: $\Hom_{\cat{C}/A}(x,y)$, per a $x\colon X\to A$ i $y\colon Y\to A$, és el conjunt dels morfismes $h\colon X\to Y$ de $\cat{C}$ que fan commutar el triangle
\[
\begin{tikzpicture}[baseline=(m.center)]
  \matrix (m) [matrix of math nodes, row sep=2.4em, column sep=2.4em] {
    X & & Y \\
      & A & \\
  };
  \draw[-{Stealth[scale=1.0]}] (m-1-1) -- node[above] {$h$} (m-1-3);
  \draw[-{Stealth[scale=1.0]}] (m-1-1) -- node[below left] {$x$} (m-2-2);
  \draw[-{Stealth[scale=1.0]}] (m-1-3) -- node[below right] {$y$} (m-2-2);
\end{tikzpicture}
\qquad\text{és a dir, } y\comp h=x;
\]
\item composició: la de $\cat{C}$; si $h\in\Hom_{\cat{C}/A}(x,y)$ i $k\in\Hom_{\cat{C}/A}(y,z)$, aleshores $k\comp h\in\Hom_{\cat{C}/A}(x,z)$, perquè $z\comp(k\comp h)=(z\comp k)\comp h=y\comp h=x$;
\item identitats: la identitat de $x\colon X\to A$ és $\id_X$, que fa commutar el triangle perquè $x\comp\id_X=x$.
\end{itemize}
Els axiomes s'hereten de $\cat{C}$.
\end{definicio}

Anàlogament es defineix la \textbf{categoria de $\cat{C}$ sota $A$} (en anglès, \emph{coslice category}), que escrivim $A/\cat{C}$:
\begin{itemize}
\item objectes: $\Ob(A/\cat{C})$ és la col·lecció dels morfismes $x\colon A\to X$ de $\cat{C}$ amb domini $A$;
\item morfismes: $\Hom_{A/\cat{C}}(x,y)$, per a $x\colon A\to X$ i $y\colon A\to Y$, és el conjunt dels morfismes $h\colon X\to Y$ de $\cat{C}$ amb $h\comp x=y$;
\item composició i identitats: les de $\cat{C}$, com a $\cat{C}/A$.
\end{itemize}
Un morfisme de $A/\cat{C}$ és, doncs, un $h$ que fa commutar el triangle sota $A$:
\[
\begin{tikzpicture}[baseline=(m.center)]
  \matrix (m) [matrix of math nodes, row sep=2.4em, column sep=2.4em] {
      & A & \\
    X & & Y \\
  };
  \draw[-{Stealth[scale=1.0]}] (m-1-2) -- node[above left] {$x$} (m-2-1);
  \draw[-{Stealth[scale=1.0]}] (m-1-2) -- node[above right] {$y$} (m-2-3);
  \draw[-{Stealth[scale=1.0]}] (m-2-1) -- node[below] {$h$} (m-2-3);
\end{tikzpicture}
\]
Aquí també convé fer explícits els extrems, com a l'observació~\ref{obs:extrems-implicits}. Un morfisme de $\cat{C}/A$ de $x$ a $y$ no és el morfisme $h\colon X\to Y$ de $\cat{C}$ tot sol, sinó $h$ \emph{junt amb} els seus extrems $x$ i $y$ a $\cat{C}/A$: la terna $(x,y,h)$. La raó és la mateixa: un mateix $h\colon X\to Y$ de $\cat{C}$ pot fer commutar el triangle per a objectes distints de $\cat{C}/A$ ---per exemple, si $X$ és el domini de dos objectes diferents $x,x'\colon X\to A$---, i sense registrar els extrems tindria dominis i codominis indeterminats a $\cat{C}/A$. Escriurem $h$ quan els extrems se sobreentenguin. Anàlogament per a $A/\cat{C}$.

\begin{exemple}
A $\cat{Set}/A$, un objecte és una aplicació $x\colon X\to A$. Les fibres $x^{-1}(a)$, per a $a\in A$, permeten interpretar aquest objecte com una família de conjunts $\bigl(x^{-1}(a)\bigr)_{a\in A}$ indexada per $A$. Recíprocament, una família de conjunts indexada per $A$ dona, mitjançant la unió disjunta, una aplicació cap a $A$. Aquestes dues construccions no identifiquen literalment les dades ---la unió disjunta demana triar-ne una representació---, però sí que estableixen la correspondència adequada entre totes dues descripcions. La farem precisa als capítols~\ref{cap:functors} i~\ref{cap:transformacions}, quan disposem del llenguatge dels functors i de les equivalències de categories.
\end{exemple}

\begin{exemple}\index{conjunt apuntat}
Prenem $A=\{\ast\}$, un conjunt d'un sol element. Un objecte de $\{\ast\}/\cat{Set}$ és una aplicació $\{\ast\}\to X$, és a dir, l'elecció d'un element distingit de $X$, el \textbf{punt base}. Un morfisme és una aplicació que respecta el punt base. Així, $\{\ast\}/\cat{Set}$ és la categoria dels \textbf{conjunts apuntats} (en anglès, \emph{pointed sets}).
\end{exemple}

\section{Principi de dualitat}\label{sec:dualitat}\index{dualitat!principi de}

La dualitat no és només una llista de parelles de conceptes. És un principi que transforma definicions, enunciats i demostracions, i la categoria oposada permet formular-lo amb precisió. Aquesta secció en fa explícit el mecanisme; el retrobarem en tractar objectes inicials i terminals, productes i coproductes, productes fibrats i coproductes fibrats (en anglès, \emph{pullbacks} i \emph{pushouts}), i límits i colímits.

\subsection{Dualitzar un enunciat}\index{dualitat!d'un enunciat}

El principi s'aplica a afirmacions escrites \emph{exclusivament} en el llenguatge de les categories: hi poden intervenir objectes i morfismes, el domini i el codomini de cada morfisme, identitats, composicions i igualtats entre composicions, combinats amb connectors lògics i amb quantificadors sobre objectes i morfismes. Les definicions de monomorfisme, d'isomorfisme o de secció són d'aquesta mena. En canvi, \enquote{$f$ és injectiva} no ho és, perquè parla dels elements d'un conjunt.

Donada una afirmació $\varphi$ d'aquest tipus, n'obtenim la \textbf{dual} $\varphi\op$ \emph{invertint totes les fletxes}:
\begin{itemize}
  \item cada morfisme $f\colon A\to B$ s'inverteix: a la dual es llegeix com una fletxa $B\to A$; per tant, s'intercanvien domini i codomini;
  \item cada composició $g\comp f$ passa a ser $f\comp g$, és a dir, s'inverteix l'ordre de composició;
  \item les identitats, les igualtats, els connectors i els quantificadors es mantenen.
\end{itemize}
Dualitzar dues vegades retorna l'afirmació de partida: $(\varphi\op)\op=\varphi$.

La dual té una lectura immediata: $\varphi\op$ diu de $\cat{C}$ exactament el que $\varphi$ diu de $\cat{C}\op$. En efecte, un morfisme $A\to B$ de $\cat{C}\op$ és un morfisme $B\to A$ de $\cat{C}$, les identitats són les mateixes i la composició de $\cat{C}\op$ és la de $\cat{C}$ en l'ordre invers, $g\op\comp f\op=(f\comp g)\op$. Per tant, per a tota categoria $\cat{C}$,
\[
\varphi \text{ es compleix a } \cat{C}\op
\quad\Longleftrightarrow\quad
\varphi\op \text{ es compleix a } \cat{C}.
\]
D'aquí surt el principi.

\begin{teorema}[Principi de dualitat]\label{teo:dualitat}\index{dualitat!principi de}
Sigui $\varphi$ una afirmació expressable en el llenguatge de les categories. Si $\varphi$ és vàlida en totes les categories, la seva dual $\varphi\op$ també és vàlida en totes les categories.
\end{teorema}

\begin{prova}
Sigui $\cat{C}$ una categoria qualsevol. Com que $\cat{C}\op$ també és una categoria, per hipòtesi $\varphi$ es compleix a $\cat{C}\op$; per l'equivalència anterior, $\varphi\op$ es compleix a $\cat{C}$.
\end{prova}

La manera habitual d'usar el principi és aquesta: s'aplica un teorema conegut a $\cat{C}\op$ i se'n tradueix la conclusió al llenguatge de $\cat{C}$. No cal repetir la demostració; \textbf{la demostració dual és la demostració original amb totes les fletxes invertides}.

Aquesta formulació operativa és la que farem servir al llarg del llibre. Se'n pot donar també una versió completament formal: fixar un llenguatge de primer ordre per a la teoria de categories, definir la dualització com una transformació de fórmules i comprovar que els axiomes en queden invariants. Aquesta formalització, amb una versió sintàctica (sobre deduccions) i una de semàntica (sobre models), és a l'apèndix~\ref{cap:axiomatica-relacional} (secció~\ref{sec:ar-dualitat}); no és necessària per seguir la resta del llibre.

\subsection{Primer exemple: de monomorfisme a epimorfisme}

Recordem (subsecció~\ref{subsec:mono-epi}) que $m\colon A\to B$ és un \textbf{monomorfisme} si, per a tot parell $u,v\colon X\to A$,
\[
\begin{tikzpicture}[baseline=(m.center)]
  \matrix (m) [matrix of math nodes, column sep=3.2em] { X & A & B \\ };
  \draw[-{Stealth[scale=1.0]}] ([yshift=2.6pt]m-1-1.east) -- node[above] {$u$} ([yshift=2.6pt]m-1-2.west);
  \draw[-{Stealth[scale=1.0]}] ([yshift=-2.6pt]m-1-1.east) -- node[below] {$v$} ([yshift=-2.6pt]m-1-2.west);
  \draw[-{Stealth[scale=1.1]}] (m-1-2) -- node[above] {$m$} (m-1-3);
\end{tikzpicture}
\qquad
m\comp u = m\comp v \implies u=v.
\]
Formem-ne el dual: intercanviem domini i codomini i invertim l'ordre de cada composició. L'enunciat dual diu que, per a tot parell $u,v\colon A\to X$,
\[
\begin{tikzpicture}[baseline=(m.center)]
  \matrix (m) [matrix of math nodes, column sep=3.2em] { B & A & X \\ };
  \draw[-{Stealth[scale=1.1]}] (m-1-1) -- node[above] {$m$} (m-1-2);
  \draw[-{Stealth[scale=1.0]}] ([yshift=2.6pt]m-1-2.east) -- node[above] {$u$} ([yshift=2.6pt]m-1-3.west);
  \draw[-{Stealth[scale=1.0]}] ([yshift=-2.6pt]m-1-2.east) -- node[below] {$v$} ([yshift=-2.6pt]m-1-3.west);
\end{tikzpicture}
\qquad
u\comp m = v\comp m \implies u=v,
\]
que és exactament la definició d'\textbf{epimorfisme}. El dual de \enquote{ser un monomorfisme} és, doncs, \enquote{ser un epimorfisme}.

Ara podem dualitzar teoremes sense repetir-ne la demostració. Per exemple, en qualsevol categoria, si $A\xrightarrow{f}B\xrightarrow{g}C$ són monomorfismes, la composició $g\comp f$ també ho és. Pel principi de dualitat (teorema~\ref{teo:dualitat}), l'enunciat dual val igualment:
\begin{quote}
en qualsevol categoria, la composició de dos epimorfismes és un epimorfisme.
\end{quote}
No n'hem hagut de fer cap demostració nova. L'exercici~\ref{exr:pr-dual-mono-epi} de la Pràctica fa aquest pas en detall i assenyala exactament què s'inverteix en passar a $\cat{C}\op$.

\subsection{Segon exemple: de secció a retracció}

Recordem (subsecció~\ref{subsec:seccions-retraccions}) que $s\colon B\to A$ és una \textbf{secció} de $f\colon A\to B$ quan $f\comp s=\id_B$. Dualitzem aquest enunciat: $f$ passa a ser una fletxa $B\to A$, $s$ una fletxa $A\to B$, i la igualtat $f\comp s=\id_B$ es converteix en $s\comp f=\id_B$. En diagrames, invertim les tres fletxes del triangle:
\[
\begin{tikzpicture}[baseline=(m.center)]
  \matrix (m) [matrix of math nodes, row sep=2.6em, column sep=2.6em] {
    B & A \\
    B & \\
  };
  \draw[-{Stealth[scale=1.1]}] (m-1-1) -- node[above] {$s$} (m-1-2);
  \draw[-{Stealth[scale=1.1]}] (m-1-2) -- node[right] {$f$} (m-2-1);
  \draw[-{Stealth[scale=1.1]}] (m-1-1) -- node[left] {$\id_B$} (m-2-1);
\end{tikzpicture}
\quad f\comp s=\id_B
\qquad\qquad
\begin{tikzpicture}[baseline=(m.center)]
  \matrix (m) [matrix of math nodes, row sep=2.6em, column sep=2.6em] {
    B & A \\
    B & \\
  };
  \draw[-{Stealth[scale=1.1]}] (m-1-2) -- node[above] {$s$} (m-1-1);
  \draw[-{Stealth[scale=1.1]}] (m-2-1) -- node[right] {$f$} (m-1-2);
  \draw[-{Stealth[scale=1.1]}] (m-2-1) -- node[left] {$\id_B$} (m-1-1);
\end{tikzpicture}
\quad s\comp f=\id_B
\]
La igualtat de la dreta diu exactament que $s$ és una \textbf{retracció} de $f$. El dual de \enquote{tenir una secció} és, doncs, \enquote{tenir una retracció}, i el d'epimorfisme escindit és monomorfisme escindit.

Amb aquest diccionari es poden dualitzar també els resultats sobre seccions i retraccions; els exercicis proposats del capítol ho apliquen a la proposició~\ref{prop:escindit-mono-epi}.

\subsection{Tercer exemple: del producte al coproducte}\label{subsec:producte-coproducte}

A tall d'avançament ---els definirem en rigor al capítol de límits---, considerem el producte i el coproducte de dos objectes, que il·lustren la dualització d'un enunciat amb un quantificador d'existència única.

Un \emph{producte} de $A$ i $B$ és un objecte $A\times B$ amb dos morfismes $\pi_1\colon A\times B\to A$ i $\pi_2\colon A\times B\to B$ tals que, per a tot $X$ i tot parell $f\colon X\to A$, $g\colon X\to B$, existeix un únic $u\colon X\to A\times B$ amb $\pi_1\comp u=f$ i $\pi_2\comp u=g$:
\[
\begin{tikzpicture}[baseline=(m.center)]
  \matrix (m) [matrix of math nodes, row sep=2.6em, column sep=2.8em] {
     & X & \\
    A & A\times B & B \\
  };
  \draw[-{Stealth[scale=1.0]}] (m-1-2) -- node[above left] {$f$} (m-2-1);
  \draw[-{Stealth[scale=1.0]}] (m-1-2) -- node[above right] {$g$} (m-2-3);
  \draw[-{Stealth[scale=1.0]},dashed] (m-1-2) -- node[right] {$u$} (m-2-2);
  \draw[-{Stealth[scale=1.0]}] (m-2-2) -- node[below] {$\pi_1$} (m-2-1);
  \draw[-{Stealth[scale=1.0]}] (m-2-2) -- node[below] {$\pi_2$} (m-2-3);
\end{tikzpicture}
\]
En el diagrama seguim una convenció que farem servir cada cop que dibuixem una propietat universal: la \textbf{fletxa discontínua} marca el morfisme l'existència del qual afirma la propietat ---aquí el mediador $u$, un cop donades la resta de dades: l'objecte $X$ i les fletxes $f$ i $g$---, mentre que les fletxes contínues són les dades de partida. Igualment, escriurem $\exists!$ per abreujar \enquote{existeix un únic}, com acabem de fer en enunciar la propietat.

Invertint totes les fletxes obtenim un objecte $A+B$ amb dos morfismes $\iota_1\colon A\to A+B$ i $\iota_2\colon B\to A+B$ tals que, per a tot $X$ i tot parell $f\colon A\to X$, $g\colon B\to X$, existeix un únic $u\colon A+B\to X$ amb $u\comp\iota_1=f$ i $u\comp\iota_2=g$:
\[
\begin{tikzpicture}[baseline=(m.center)]
  \matrix (m) [matrix of math nodes, row sep=2.6em, column sep=2.8em] {
    A & A+B & B \\
     & X & \\
  };
  \draw[-{Stealth[scale=1.0]}] (m-1-1) -- node[above] {$\iota_1$} (m-1-2);
  \draw[-{Stealth[scale=1.0]}] (m-1-3) -- node[above] {$\iota_2$} (m-1-2);
  \draw[-{Stealth[scale=1.0]}] (m-1-1) -- node[below left] {$f$} (m-2-2);
  \draw[-{Stealth[scale=1.0]}] (m-1-3) -- node[below right] {$g$} (m-2-2);
  \draw[-{Stealth[scale=1.0]},dashed] (m-1-2) -- node[right] {$u$} (m-2-2);
\end{tikzpicture}
\]
Aquesta és la definició del \emph{coproducte} $A+B$. Els quantificadors $\forall$ i $\exists!$ i la conjunció de les dues equacions es mantenen; el que s'inverteix és el domini i el codomini de cada morfisme i l'ordre de cada composició.

\paragraph{Un cas que cal vigilar.} La construcció \enquote{categoria sobre $A$} es dualitza en la construcció \enquote{categoria sota $A$} (definició~\ref{def:slice}). Ara bé, en dualitzar aquesta construcció cal dualitzar també la categoria ambient. La fórmula precisa és
\[
(\cat{C}/A)\op\;\cong\;A/(\cat{C}\op),
\qquad
(A/\cat{C})\op\;\cong\;(\cat{C}\op)/A.
\]
No s'ha de llegir, doncs, que $(\cat{C}/A)\op\cong A/\cat{C}$ mantenint $\cat{C}$ intacta: això és fals en general. Per exemple, sigui $\cat{C}=\cat{2}=(0\to 1)$ i prenguem com a $A$ l'\emph{objecte} $1$ de $\cat{2}$ (no la categoria $\cat{1}$). La categoria $\cat{2}/1$ té dos objectes (les dues fletxes amb codomini $1$: la fletxa $0\to 1$ i $\id_1$), mentre que $1/\cat{2}$ en té un de sol ($\id_1$, l'única fletxa amb domini $1$); per tant $(\cat{2}/1)\op$ \emph{no} és isomorfa a $1/\cat{2}$. El que sí que val és $(\cat{2}/1)\op\cong 1/(\cat{2}\op)$: a $\cat{2}\op$ hi ha dues fletxes amb domini $1$, la identitat i la fletxa $1\to 0$. En tots els casos, la mateixa regla ---invertir les fletxes--- transforma cada noció en la seva dual, però la regla actua sobre \emph{tot} l'enunciat, ambient inclòs.

\subsection{Diccionari bàsic de dualitats}

Recollim les parelles duals. Les primeres ja han aparegut en aquest capítol; la resta les trobarem més endavant, i la taula creixerà a mesura que apareguin nous conceptes.
\begin{center}
\begin{tabular}{ll}
\hline
\textbf{Concepte} & \textbf{Concepte dual}\\
\hline
\multicolumn{2}{l}{\emph{Ja vistos en aquest capítol}}\\
monomorfisme & epimorfisme\\
secció & retracció\\
producte & coproducte\\
categoria sobre $A$ \ ($\cat{C}/A$) & categoria sota $A$ \ ($A/\cat{C}$)\\
\hline
\multicolumn{2}{l}{\emph{Més endavant}}\\
objecte inicial & objecte terminal\\
equalitzador & coequalitzador\\
producte fibrat & coproducte fibrat\\
con & cocon\\
límit & colímit\\
\hline
\end{tabular}
\end{center}
Algunes nocions són \textbf{autoduals}\index{dualitat!noció autodual}: la identitat, l'isomorfisme i la commutativitat d'un diagrama continuen sent nocions del mateix tipus després de dualitzar-les.

La taula recull les parelles duals com a \emph{nocions}: cada concepte i el seu dual. Quan aquestes nocions es realitzen en categories concretes, però, la dualitat s'ha d'aplicar a tot l'enunciat, ambient inclòs. La parella \enquote{sobre $A$ / sota $A$} n'és l'exemple a vigilar: com hem vist, l'enunciat correcte és $(\cat{C}/A)\op\cong A/(\cat{C}\op)$, no $(\cat{C}/A)\op\cong A/\cat{C}$ amb $\cat{C}$ intacta.

\subsection{Abast i precaucions}

El principi només s'aplica a afirmacions expressables en el llenguatge de les categories. Una afirmació que faci servir informació externa ---elements d'un conjunt, oberts d'un espai o una construcció concreta--- no té dual fins que no es reformula exclusivament amb objectes i morfismes. I quan l'afirmació es refereix a construccions fetes a partir d'una categoria, cal dualitzar també la categoria ambient, com acabem de veure amb les categories sobre i sota un objecte.

El principi de dualitat no permet concloure, per si sol, que una categoria concreta sigui equivalent a la seva oposada. El principi relaciona afirmacions sobre $\cat{C}$ amb les afirmacions duals sobre $\cat{C}\op$; no proporciona en general cap equivalència $\cat{C}\simeq\cat{C}\op$. Algunes categories sí que ho són ---per exemple, $\cat{2}=(0\to 1)$ és isomorfa a $\cat{2}\op$ intercanviant-ne els dos objectes---, però això s'ha de comprovar en cada cas. En canvi, $\cat{Set}$ no s'identifica amb $\cat{Set}\op$, encara que els teoremes generals sobre productes tinguin teoremes duals sobre coproductes.

Finalment, una precaució sobre la lògica. Els connectors i els quantificadors amb què escrivim una afirmació $\varphi$ pertanyen al llenguatge amb què \emph{parlem} de la categoria, i la dualitat els deixa intactes. Una altra cosa és el que passa quan els objectes de la categoria són ells mateixos fórmules, com a la categoria $\cat{Prov}_T$ de la subsecció~\ref{subsec:elementals}. Allà, \emph{quan existeixen}, la conjunció és un producte, la disjunció un coproducte, el cert un objecte terminal i el fals un objecte inicial, i la dualitat de fletxes els intercanvia. Però un preordre de deduïbilitat qualsevol no garanteix que aquests connectors existeixin: en el preordre discret de dues fórmules sense cap deducció entre elles, aquestes dues fórmules no tenen ni producte ni coproducte. I que la dualitat de fletxes coincideixi exactament amb la dualitat de De~Morgan exigeix, a més, un complement involutiu com el de les àlgebres booleanes. Tornarem sobre aquesta lògica interna més endavant, amb el llenguatge adequat.

\section{Propietats duals dels morfismes}\label{sec:propietats-duals}

El principi de dualitat estalvia feina: n'hi ha prou de demostrar les propietats dels monomorfismes; les dels epimorfismes se n'obtenen invertint les fletxes, sense cap demostració nova. Ho il·lustrem amb les propietats bàsiques.

\subsection{Propietats dels monomorfismes}

\begin{proposicio}\label{prop:mono-comp}
La composició de dos monomorfismes és un monomorfisme: si $f\colon A\to B$ i $g\colon B\to C$ són monomorfismes, aleshores $g\comp f$ ho és.
\end{proposicio}

\begin{prova}
Siguin $u,v\colon X\to A$ amb $(g\comp f)\comp u=(g\comp f)\comp v$, és a dir $g\comp(f\comp u)=g\comp(f\comp v)$. Com que $g$ és mono, $f\comp u=f\comp v$; i com que $f$ és mono, $u=v$.
\[
\begin{tikzpicture}[baseline=(m.center)]
  \matrix (m) [matrix of math nodes, column sep=3em] { X & A & B & C \\ };
  \draw[-{Stealth[scale=1.0]}] ([yshift=2.6pt]m-1-1.east) -- node[above] {$u$} ([yshift=2.6pt]m-1-2.west);
  \draw[-{Stealth[scale=1.0]}] ([yshift=-2.6pt]m-1-1.east) -- node[below] {$v$} ([yshift=-2.6pt]m-1-2.west);
  \draw[-{Stealth[scale=1.1]}] (m-1-2) -- node[above] {$f$} (m-1-3);
  \draw[-{Stealth[scale=1.1]}] (m-1-3) -- node[above] {$g$} (m-1-4);
\end{tikzpicture}
\]
\end{prova}

\begin{proposicio}\label{prop:mono-factor}
Si la composició $g\comp f$ és un monomorfisme, aleshores $f$ és un monomorfisme.
\end{proposicio}

\begin{prova}
Siguin $u,v\colon X\to A$ amb $f\comp u=f\comp v$. Component amb $g$, $(g\comp f)\comp u=(g\comp f)\comp v$; com que $g\comp f$ és mono, $u=v$. (No cal cap hipòtesi sobre $g$.)
\end{prova}

\begin{proposicio}\label{prop:iso-retr-mono}
Tot monomorfisme escindit és un monomorfisme, i tot isomorfisme és un monomorfisme escindit. Per tant,
\[
\text{isomorfisme} \;\Longrightarrow\; \text{monomorfisme escindit} \;\Longrightarrow\; \text{monomorfisme}.
\]
\end{proposicio}

\begin{prova}
La primera implicació és la proposició~\ref{prop:escindit-mono-epi}: si $f$ té una retracció, aleshores $f$ és un monomorfisme. Per a la segona, si $f$ és un isomorfisme, el seu invers $f^{-1}$ satisfà $f^{-1}\comp f=\id_A$, de manera que $f^{-1}$ és una retracció de $f$; per tant $f$ té retracció i és un monomorfisme escindit. La cadena d'implicacions se'n segueix.
\end{prova}

\subsection{Les propietats duals dels epimorfismes}

Aplicant el principi de dualitat a les proposicions anteriors ---invertint les fletxes--- obtenim, sense cap demostració nova, les propietats corresponents dels epimorfismes.

\begin{proposicio}[dual de la proposició~\ref{prop:mono-comp}]\label{prop:epi-comp}
La composició de dos epimorfismes és un epimorfisme.
\end{proposicio}

\begin{proposicio}[dual de la proposició~\ref{prop:mono-factor}]\label{prop:epi-factor}
Si la composició $g\comp f$ és un epimorfisme, aleshores $g$ és un epimorfisme.
\end{proposicio}

\begin{proposicio}[dual de la proposició~\ref{prop:iso-retr-mono}]\label{prop:iso-sec-epi}
Tot epimorfisme escindit és un epimorfisme, i tot isomorfisme és un epimorfisme escindit. Per tant,
\[
\text{isomorfisme} \;\Longrightarrow\; \text{epimorfisme escindit} \;\Longrightarrow\; \text{epimorfisme}.
\]
\end{proposicio}

Observem l'asimetria que revela la dualitat: a la proposició~\ref{prop:mono-factor}, d'un mono $g\comp f$ hereta el factor \emph{de la dreta}, $f$; a la seva dual~\ref{prop:epi-factor}, d'un epi hereta el factor \emph{de l'esquerra}, $g$.

\subsection{Condicions necessàries i suficients}

Les implicacions anteriors \emph{no} es reverteixen en general, i és aquí on convé anar amb compte. La cadena vàlida a qualsevol categoria és
\[
\text{iso} \;\Longrightarrow\; \text{mono escindit} \;\overset{(1)}{\Longrightarrow}\; \text{mono},
\]
i cap de les dues fletxes es pot invertir en general:
\begin{itemize}
\item[(1)] \textbf{mono $\not\Rightarrow$ mono escindit.} A $\cat{Grp}$, l'homomorfisme $\mathbb{Z}\xrightarrow{\,\cdot 2\,}\mathbb{Z}$ és un monomorfisme, però no té cap retracció: no és escindit.
\item[(2)] \textbf{mono escindit $\not\Rightarrow$ iso.} Un monomorfisme escindit no ha de ser exhaustiu en cap sentit: a $\cat{Set}$, la injecció $\{0\}\hookrightarrow\{0,1\}$ té retracció (l'única aplicació $\{0,1\}\to\{0\}$), és per tant un monomorfisme escindit, però no és un isomorfisme.
\end{itemize}
Dualment, $\text{iso}\Rightarrow\text{epi escindit}\Rightarrow\text{epi}$, amb les mateixes dues fletxes no reversibles. Convé no confondre \enquote{mono escindit $+$ epi} amb \enquote{mono $+$ epi}: un \textbf{bimorfisme} (alhora mono i epi) no és necessàriament un isomorfisme ---l'única fletxa no trivial de $\cat{2}$ ho és sense tenir invers, i a $\cat{Top}$ una bijecció contínua amb inversa no contínua també---, mentre que un morfisme alhora mono escindit i epi sí que ho és, com recull la caracterització següent:
\[
f \text{ és iso} \iff f \text{ és mono escindit i epi} \iff f \text{ és epi escindit i mono}.
\]
Les demostracions d'aquestes equivalències, i la de la cadena d'implicacions anterior, es proposen a les parts de pràctica i d'exercicis.

\section{Grafs i categories lliures}\label{sec:grafs-lliures}\index{categoria!lliure}

Els grafs i les categories estan estretament relacionats, però no són el mateix. Un graf dirigit només especifica vèrtexs i arestes. Una categoria, a més, disposa d'identitats, permet compondre fletxes i exigeix que aquestes composicions satisfacin els axiomes. De fet, per a una categoria petita, si oblidem la composició i quines fletxes són identitats, els objectes i morfismes formen un graf dirigit: és el seu \emph{graf subjacent}. Les fletxes identitat no desapareixen ---continuen sent arestes del graf, un bucle a cada vèrtex---, només que ja no estan marcades com a identitats. Cal que la categoria sigui petita, perquè els vèrtexs i les arestes d'un graf formen conjunts.

Hi ha una manera natural de fer el camí contrari: generar una categoria a partir d'un graf dirigit prenent com a morfismes tots els camins possibles.

\begin{definicio}\label{def:free}
Sigui $G$ un graf dirigit, amb conjunt de vèrtexs $V_G$, conjunt d'arestes $E_G$ i aplicacions origen i destí $s_G,t_G\colon E_G\to V_G$. Per a $A,B\in V_G$, un \textbf{camí}\index{camí} de $A$ a $B$ de longitud $n\ge 1$ és una successió $(e_1,\dots,e_n)$ d'arestes amb
\[
s_G(e_1)=A,\qquad t_G(e_i)=s_G(e_{i+1})\ \ (1\le i<n),\qquad t_G(e_n)=B.
\]
Per a cada vèrtex $A$ hi ha, a més, el \textbf{camí buit} $()_A$ de $A$ a $A$, de longitud zero.

La \textbf{categoria lliure generada pel graf} $G$, que escrivim $\Free(G)$,\index{Free@$\Free$}\index{categoria!lliure} té:
\begin{itemize}
\item objectes: $\Ob(\Free(G))=V_G$, els vèrtexs de $G$;
\item morfismes: $\Hom_{\Free(G)}(A,B)$ és el conjunt dels camins de $A$ a $B$;
\item composició: la \textbf{concatenació} en ordre de recorregut; per a $p=(a_1,\dots,a_r)\in\Hom(A,B)$ i $q=(b_1,\dots,b_s)\in\Hom(B,C)$,
\[
q\comp p=(a_1,\dots,a_r,b_1,\dots,b_s),
\]
on $p$, que surt de $A$, es recorre primer, d'acord amb la convenció $q\comp p$; concatenar amb un camí buit deixa el camí igual;
\item identitats: $\id_A=()_A$.
\end{itemize}
\end{definicio}

Per exemple, a partir del graf
\[
A\xrightarrow{f}B\xrightarrow{g}C
\]
la categoria lliure conté, a més de les identitats i dels camins $(f)$ i $(g)$ de longitud~1, que escrivim simplement $f$ i $g$, la fletxa composta $g\comp f=(f,g)\colon A\to C$:
\[
\begin{tikzpicture}[baseline=(m.center)]
  \matrix (m) [matrix of math nodes, column sep=3.5em] { A & B & C \\ };
  \draw[-{Stealth[scale=1.05]}] (m-1-1) -- node[above] {$f$} (m-1-2);
  \draw[-{Stealth[scale=1.05]}] (m-1-2) -- node[above] {$g$} (m-1-3);
  \draw[-{Stealth[scale=1.0]}] (m-1-1) to[bend right=25] node[below] {$g\comp f$} (m-1-3);
\end{tikzpicture}
\]
Els axiomes de categoria se satisfan. La composició és associativa, perquè $r\comp(q\comp p)$ i $(r\comp q)\comp p$ són totes dues la successió de totes les arestes de $p$, $q$ i $r$, en aquest ordre; i el camí buit és neutre, perquè concatenar-lo no afegeix cap aresta: $p\comp()_A=p=()_B\comp p$ per a tot camí $p$ de $A$ a $B$. Com que $V_G$ i $E_G$ són conjunts, també ho són els conjunts de camins, i $\Free(G)$ és una categoria petita. L'apèndix~\ref{cap:grafs} desenvolupa aquesta construcció amb més exemples. El nom \emph{lliure} quedarà justificat més endavant, quan disposem del llenguatge dels functors i puguem formular la propietat universal d'aquesta construcció.

\begin{exemple}[El monoide lliure i la convenció de composició]\label{ex:monoide-lliure-cat}
Prenem el cas d'un graf amb \textbf{un sol vèrtex} $\ast$ i conjunt d'arestes $E$. Aleshores totes les arestes són bucles ---comencen i acaben a $\ast$--- i cap parell no deixa d'estar encadenat, de manera que un camí és simplement una seqüència finita d'arestes, és a dir, una \emph{paraula} sobre $E$. Escrivim $E^*$ per al conjunt de totes aquestes paraules, incloent-hi la \emph{paraula buida} $\varepsilon$, de longitud zero.

Fixem tres convencions, que s'han de llegir juntes.
\begin{itemize}
\item \emph{Lectura de les paraules.} Escrivim les lletres en l'\emph{ordre de recorregut}: la paraula $a_1a_2\cdots a_n$ és el camí $(a_1,a_2,\dots,a_n)$, que recorre primer $a_1$, després $a_2$, i així successivament.
\item \emph{Concatenació.} Per a $u=a_1\cdots a_m$ i $v=b_1\cdots b_n$, la concatenació és la juxtaposició $uv=a_1\cdots a_m b_1\cdots b_n$: primer les lletres de $u$ i després les de $v$. Amb aquesta operació i la paraula buida, $(E^*,\text{concatenació},\varepsilon)$ és un monoide, el \textbf{monoide lliure}\index{monoide!lliure} sobre $E$: $(uv)w$ i $u(vw)$ són totes dues la paraula que juxtaposa $u$, $v$ i $w$ en aquest ordre, i $\varepsilon u=u=u\varepsilon$.
\item \emph{Composició.} A $\Free(G)$, la composició és la de la definició: $v\comp u$ vol dir \enquote{primer $u$, després $v$}, i el camí que en resulta recorre primer les lletres de $u$. Per tant,
\[
v\comp u=uv.
\]
\end{itemize}
Amb dues lletres $a,b\in E$, tot cap en una línia:
\[
b\comp a=ab\quad(\text{primer }a,\text{ després }b),
\qquad\qquad
a\comp b=ba\quad(\text{primer }b,\text{ després }a).
\]
L'ordre de la composició és, doncs, el \emph{contrari} del de la juxtaposició: el símbol de la dreta a $b\comp a$ és el que s'escriu primer a la paraula.

Quina categoria hem obtingut? Té un sol objecte, $\Hom_{\Free(G)}(\ast,\ast)=E^*$ i $\id_\ast=\varepsilon$. Però, amb la convenció de $\cat{B}M$, en què $g\comp f=g\cdot f$, la composició $v\comp u=uv$ correspon al producte \emph{invertit}. Si anomenem \textbf{monoide oposat}\index{monoide!oposat} $M\op$ d'un monoide $M$ el que té els mateixos elements i el mateix neutre i l'operació $m\cdot_{\mathrm{op}} n=n\cdot m$ ---de manera que $\cat{B}(M\op)=(\cat{B}M)\op$---, el que tenim és la \emph{igualtat}
\[
\Free(G)=\cat{B}\bigl((E^*)\op\bigr)=\bigl(\cat{B}(E^*)\bigr)\op,
\]
i no $\cat{B}(E^*)$. La diferència és només d'ordre d'escriptura. L'aplicació que inverteix l'ordre de les lletres, $a_1\cdots a_n\mapsto a_n\cdots a_1$, compleix $(uv)^{\mathrm{rev}}=v^{\mathrm{rev}}u^{\mathrm{rev}}$, i és per tant un isomorfisme de monoides $E^*\cong(E^*)\op$. Dona, doncs, un isomorfisme canònic de categories $\cat{B}(E^*)\cong\Free(G)$. Si haguéssim escrit les paraules en l'ordre de la composició ($a_n\cdots a_1$ per al camí que comença per $a_1$), hauríem obtingut literalment $\cat{B}(E^*)$. Hem preferit l'ordre de recorregut, que és el natural per als camins. En aquest sentit, i llevat d'aquest isomorfisme, el monoide lliure és el cas particular de la categoria lliure per a un graf amb un sol vèrtex.
\end{exemple}

\begin{exemple}
La categoria lliure sobre el graf amb dos vèrtexs $A,B$ i una sola aresta $A\to B$ és la categoria finita $\mathbf{2}$. Si, en canvi, el graf té dues arestes en sentits oposats, $e\colon A\to B$ i $f\colon B\to A$, la categoria lliure té infinits morfismes: $e$, $f$, $f\comp e$, $e\comp f$, $e\comp f\comp e$, i així successivament, perquè cap relació no identifica camins diferents. Així, \emph{un graf finit no genera necessàriament una categoria lliure finita}.
\end{exemple}

Aquesta construcció, que genera l'estructura més econòmica a partir d'unes dades sense imposar cap relació que no hi vingui forçada, és un exemple del patró \emph{lliure}, que reapareix per a monoides, grups i moltes altres estructures i que, al capítol~\refcap{cap:adjuncions}{6}, es formalitza com una \emph{adjunció} entre un functor lliure i un functor oblit.

\section{Categories segons la mida}\index{categoria!petita}\index{categoria!localment petita}
\label{sec:mida}

Hem parlat de \enquote{col·lecció} d'objectes en lloc de \enquote{conjunt}, i ara podem precisar per què. En alguns exemples, els objectes són massa nombrosos per formar un conjunt.

\begin{observacio}[La paradoxa de Russell]
La categoria $\cat{Set}$ té per objectes \emph{tots} els conjunts. Ara bé, no pot existir el conjunt de tots els conjunts. En efecte, si $U$ fos aquest conjunt, en podríem separar el subconjunt $R=\{x\in U\mid x\notin x\}$ dels conjunts que no es pertanyen a si mateixos, i preguntar-nos si $R\in R$: si $R\in R$, la definició força $R\notin R$; i si $R\notin R$, la definició força $R\in R$. Totes dues possibilitats són contradictòries. Aquesta és la \textbf{paradoxa de Russell}, i mostra que la col·lecció de tots els conjunts és massa gran per ser un conjunt.
\end{observacio}

Per això, en parlar dels objectes d'una categoria, diem \enquote{col·lecció} i no \enquote{conjunt}: pot ser massa gran. Segons la mida d'aquestes col·leccions introduïm els conceptes següents.

\begin{definicio}
Una categoria és \textbf{petita} si la seva col·lecció d'objectes i la seva col·lecció de morfismes són tots dos conjunts. En cas contrari es diu \textbf{gran}\index{categoria!gran}. I és \textbf{localment petita} si, per a cada parella d'objectes $A,B$, la col·lecció de morfismes $\Hom(A,B)$ és un conjunt, encara que la col·lecció de tots els objectes pugui ser massa gran per ser-ho.
\end{definicio}

Convé no llegir aquestes tres nocions com una partició en tres classes disjuntes. La partició real és \emph{petita} contra \emph{gran}. La petitesa local és una condició transversal: tota categoria petita és localment petita ---si tots els morfismes formen un conjunt, també en formen un els de cada parella $(A,B)$---, i una categoria gran pot ser localment petita o no. Així,
\[
\text{petita}\;\Longrightarrow\;\text{localment petita},
\]
i el recíproc és fals: $\cat{Set}$ és localment petita però no petita. Que no és petita ho mostra la paradoxa de Russell; que és localment petita es comprova identificant cada aplicació amb el seu graf (exercici~\ref{exr:pr-set-localment-petita} de la Pràctica). Les tres possibilitats efectives són, doncs, petita (per tant localment petita), gran localment petita i gran no localment petita.

Moltes de les categories habituals que utilitzem ---$\cat{Set}$, $\cat{Grp}$, $\cat{Top}$, $\cat{Graph}$ i altres--- són localment petites però no petites: tenen massa objectes per formar un conjunt, però entre dos objectes qualssevol hi ha només un conjunt de morfismes. En canvi, cada preordre concret sobre un conjunt i cada categoria finita concreta són categories petites. És important no confondre aquestes categories petites amb les categories grans que les apleguen, com $\cat{Pos}$, la categoria de tots els ordres parcials, o la categoria de totes les categories finites: en tots dos casos hi ha una classe pròpia de conjunts subjacents. Per la mateixa raó és gran $\cat{FinSet}$, tot i que cadascun dels seus objectes és un conjunt finit. Al llarg del llibre suposarem, tret que diguem el contrari, que les categories amb què treballem són localment petites; això garanteix que sempre podem parlar del conjunt $\Hom(A,B)$, cosa que serà essencial més endavant.\footnote{Aquí prenem com a base la teoria de conjunts clàssica, amb classes pròpies; l'apèndix~\ref{cap:mida} fixa aquest marc, i \cite[§1.8]{awodey}, \cite[§1.1]{riehl}, \cite[§3.2]{leinster} i \cite[cap.~1]{perrone} en donen presentacions introductòries. També és possible prendre el mateix llenguatge estructural com a fonament, sense una teoria de conjunts prèvia: l'apèndix~\ref{cap:axiomatica-relacional} n'ofereix un tast, i \cite[cap.~3]{leinster} en fa una introducció.}

\begin{resumcap}
\apartat{Idees centrals}
\begin{enumerate}
  \item Una categoria organitza objectes i morfismes mitjançant identitats i una composició associativa. El que importa no és què són els objectes, sinó com es relacionen a través de les fletxes. Els diagrames commutatius representen igualtats entre composicions.
  \item Una mateixa forma categòrica apareix en contextos molt diferents: conjunts i aplicacions, preordres, monoides, grafs, relacions i sistemes deductius.
  \item Els isomorfismes expressen quan dos objectes tenen la mateixa estructura dins la categoria; sovint, la noció estructuralment rellevant és l'isomorfisme més que no pas la igualtat estricta. Monomorfismes i epimorfismes es defineixen per cancel·lació. A $\cat{Set}$, i en algunes categories concretes, coincideixen amb les aplicacions injectives i exhaustives, però aquesta coincidència no és vàlida en general: en una categoria prima, per exemple, tot morfisme és mono i epi.
  \item Seccions i retraccions són inverses unilaterals. Un morfisme és un monomorfisme escindit si i només si té una retracció, i és un epimorfisme escindit si i només si té una secció. Tot monomorfisme escindit és mono i tot epimorfisme escindit és epi, però les implicacions inverses no valen en general:
  \[
  \text{iso}\Longrightarrow\text{mono escindit}\Longrightarrow\text{mono},
  \qquad
  \text{iso}\Longrightarrow\text{epi escindit}\Longrightarrow\text{epi}.
  \]
  \item A partir d'una categoria en podem construir d'altres: l'oposada, productes de categories, subcategories i categories sobre i sota un objecte.
  \item El \textbf{principi de dualitat} ---invertir totes les fletxes--- converteix sistemàticament cada definició i cada resultat en el seu dual, amb la composició en l'ordre invers.
  \item Un graf dirigit no té, per si sol, composició ni identitats; la categoria lliure hi afegeix formalment tots els camins i les identitats necessàries.
  \item Les qüestions de mida obliguen a distingir categories petites de categories localment petites. Tota categoria petita és localment petita, però moltes categories habituals, com $\cat{Set}$, són localment petites sense ser petites.
\end{enumerate}
\apartat{Errors típics} Creure que \enquote{mono $+$ epi $\Rightarrow$ iso} (fals a $\cat{Top}$ i a $\cat{2}$: un bimorfisme no és sempre un isomorfisme); confondre un monomorfisme amb un monomorfisme escindit; confondre la igualtat d'objectes amb l'isomorfisme, encara que sigui canònic; aplicar els adjectius \enquote{injectiu} i \enquote{exhaustiu} a morfismes d'una categoria general, en lloc de reservar-los per a aplicacions; raonar amb elements en una categoria on els objectes no en tenen o on no separen les fletxes; escriure la composició en l'ordre equivocat en dualitzar; i oblidar la hipòtesi de petitesa local en parlar d'homconjunts.
\apartat{Lectura recomanada} \cite[cap.~1, §2.1 i §3.1]{awodey} per a la introducció paral·lela, els monomorfismes i epimorfismes i el principi de dualitat; \cite[§1.1]{leinster} per a una presentació breu i moderna; \cite[cap.~I i II.1--II.3, II.7]{maclane} per a la formulació clàssica, inclosos la dualitat, les categories producte i les categories lliures; \cite[§1.3]{barrwells} per als grafs, en una presentació orientada a la informàtica; \cite[parts~I--II, sessions~2--10]{lawvereschanuel} per a una introducció encara més elemental als morfismes, els isomorfismes i les seccions i retraccions; i \cite[§1.1.5 i §1.2]{perrone-notes} per a contraexemples elementals d'isomorfismes, monomorfismes i epimorfismes, escindits o no. Per a la presentació relacional de les categories petites, l'apèndix~\ref{cap:axiomatica-relacional}.
\end{resumcap}
